% Canon X-07 -- Applications en assembleur
% Synthèse du livre « Applications en assembleur dans l'univers du Canon X-07 »
% (A. Tonic, E. Arevian, P. Millet, J.-J. Rousseau, Éditions Neptune, 1986),
% scan x-07_AppsASM.pdf.
% Compilation : pdflatex X07_AppsASM.tex  (deux passes pour la table des matières)

\documentclass[a4paper,10pt]{article}

\usepackage[utf8]{inputenc}
\usepackage[T1]{fontenc}
\usepackage{lmodern}
\usepackage[french]{babel}
\usepackage[margin=2cm]{geometry}
\usepackage{array}
\usepackage{booktabs}
\usepackage{longtable}
\usepackage{enumitem}
\usepackage[table]{xcolor}
\usepackage[hidelinks]{hyperref}

\setlength{\parindent}{0pt}
\setlength{\parskip}{4pt}
\setlist{nosep}

\newcommand{\hex}[1]{\texttt{#1}}
\newcommand{\reg}[1]{\textsf{#1}}
\newcommand{\pg}[1]{\texorpdfstring{\textcolor{gray}{\small(p.~#1)}}{}}
% Remarque du rédacteur : coquille, incohérence ou bogue relevé dans le livre
\newcommand{\nota}[1]{\textcolor{red!60!black}{\small\textbf{N.d.R.}~#1}}

% Fiche d'identité d'un programme
\newcommand{\fiche}[3]{%
  \begin{center}\small
  \begin{tabular}{@{}ll@{}}
  \toprule
  Implantation & #1 \\
  Longueur & #2 \\
  Lancement & #3 \\
  \bottomrule
  \end{tabular}
  \end{center}}

% Listing désassemblé : adresse, octets, étiquette, instruction, commentaire
\newenvironment{asm}
  {\small\begin{longtable}{@{}>{\ttfamily}l>{\ttfamily}l>{\ttfamily}l>{\ttfamily}l>{\raggedright\arraybackslash}p{6.2cm}@{}}
   \toprule
   \normalfont\bfseries Adr. & \normalfont\bfseries Octets & \normalfont\bfseries Étiq. & \normalfont\bfseries Instruction & \textbf{Commentaire}\\
   \midrule\endhead
   \bottomrule\endfoot}
  {\end{longtable}}

% Listing source sans adresses : étiquette, instruction, commentaire
\newenvironment{src}
  {\small\begin{longtable}{@{}>{\ttfamily}l>{\ttfamily}l>{\raggedright\arraybackslash}p{8.6cm}@{}}
   \toprule
   \normalfont\bfseries Étiq. & \normalfont\bfseries Instruction & \textbf{Commentaire}\\
   \midrule\endhead
   \bottomrule\endfoot}
  {\end{longtable}}

\title{\textbf{Canon X-07 -- Applications en assembleur}\\[4pt]
\large Synthèse du livre\\
\emph{Applications en assembleur dans l'univers du Canon X-07}\\[2pt]
\normalsize A.~Tonic, E.~Arevian, P.~Millet, J.-J.~Rousseau -- Éditions Neptune, 1986}
\author{}
\date{}

\begin{document}
\maketitle

\begin{center}
\begin{minipage}{0.85\textwidth}\small
Ce livre fait suite aux \emph{Mystères du X-07}, dont il reprend les
adresses et les routines. Il comporte trois parties : « Soft » (utiliser des
routines en langage machine depuis le BASIC), « Hard » (interruptions,
entrées/sorties, interface vidéo X-720) et « Logiciel » (douze utilitaires
et jeux). Ce document reprend l'ensemble sous forme de tableaux et de
listings. Les numéros de page \pg{xx} renvoient au livre. Sauf mention
contraire, les valeurs sont en hexadécimal.

Les programmes courts ont été vérifiés en désassemblant les octets de leurs
lignes \texttt{DATA}. Les coquilles et les bogues relevés sont signalés par
\nota{}. Les longs vidages hexadécimaux des logiciels de la 3\textsuperscript{e}
partie ne sont pas recopiés : leur transcription depuis le scan serait trop
peu fiable. On renvoie aux pages du livre.
\end{minipage}
\end{center}

\tableofcontents
\newpage

%======================================================================
\section*{Présentation du livre \pg{1-10}}
\addcontentsline{toc}{section}{Présentation du livre}

\begin{center}\small
\begin{tabular}{@{}l>{\raggedright\arraybackslash}p{11cm}@{}}
\toprule
Auteur & Présentation (p.~2) \\
\midrule
André Tonic & 19 ans, études d'économie ; président du Club C7 et directeur de la publication des Éditions Neptune \\
Eddy Arevian & 23 ans, électronicien ; conseiller technique du Club C7 \\
Pierre Millet & 23 ans ; trésorier du Club C7 \\
Jean-Jacques Rousseau & 43 ans ; maître de conférences à l'université du Maine \\
\bottomrule
\end{tabular}
\end{center}

La préface est un court récit : un technicien de Kourou dépanne le lancement
d'une fusée grâce à son X-07. L'introduction présente le livre comme la suite
des \emph{Mystères du X-07} : les notions de base de l'assembleur Z80 et les
adresses de la ROM y sont supposées connues.

%======================================================================
\part*{1\textsuperscript{re} partie : Soft \pg{13-45}}
\addcontentsline{toc}{part}{1re partie : Soft}

%----------------------------------------------------------------------
\section{Utilisation de routines en langage machine \pg{15-25}}

Une routine en langage machine se lance depuis le BASIC de deux façons :
\begin{itemize}
  \item \texttt{EXEC adresse}, sans argument ;
  \item \texttt{Z=USR(adresse,argument)}, avec un argument et une valeur de
        retour.
\end{itemize}

\subsection{Où implanter la routine ? \pg{15-19}}

Les pointeurs de début et de fin de la zone BASIC sont recalculés à chaque
mise sous tension (sauf après un \texttt{SLEEP}). Il faut donc réserver la
place de la routine. Les adresses dépendent de la quantité de mémoire
(8, 12, 16, 20 ou 24\,Ko).

\begin{center}\small
\begin{tabular}{@{}>{\raggedright}p{3.3cm}>{\raggedright}p{4.5cm}>{\raggedright}p{3.6cm}>{\raggedright\arraybackslash}p{3.6cm}@{}}
\toprule
Méthode & Principe & Avantages & Inconvénients \\
\midrule
\texttt{CLEAR} & \texttt{CLEAR 50,\&H4000} fixe la limite haute de la zone BASIC (MEMSIZ, \hex{01DF}) ; les octets situés au-delà sont libres & classique, adresse fixe & à refaire à chaque mise sous tension \\
\addlinespace
Fond des fichiers RAM & la routine est placée à la fin de la zone des fichiers RAM (RAMSTR, \hex{0210} ; taille réglée par \texttt{FSET}) & pas de \texttt{CLEAR} à refaire & vérifier la place libre avec \texttt{DIR} ; ne rien écraser \\
\addlinespace
Ligne \texttt{REM} & une ligne \texttt{REM} suivie d'autant de caractères que la routine, remplacés par \texttt{POKE} & tous les codes admis ; adresse absolue ; code résident & 80 codes au plus (LCD) ; certains codes troublent l'interpréteur ; ligne impossible à éditer ou à lister \\
\addlinespace
Chaîne dans le programme & \texttt{A\$="TURLUTUTU"} : le code remplace les caractères & adresse absolue ; code résident & codes \hex{00} et \hex{22} (guillemet) interdits ; listage imprévisible \\
\addlinespace
Zone des chaînes & \texttt{A\$=STRING\$(100,"*")} & tous les codes admis ; listage normal & relogeable uniquement ; adresse à recalculer avant chaque appel ; code à réimplanter à chaque exécution \\
\addlinespace
Tableau d'entiers & le code est rangé dans un tableau \texttt{U\%()}, deux octets par élément & passage facile de plusieurs arguments & octets croisés deux à deux ; relogeable uniquement ; mise en œuvre délicate \\
\bottomrule
\end{tabular}
\end{center}

\paragraph{Calcul des adresses.}
\begin{itemize}
  \item Ligne \texttt{REM} $n$ : après \texttt{RESTORE} $n$, le pointeur des
        DATA (\hex{0328}--\hex{0329}) désigne la ligne ; le premier caractère
        après \texttt{REM} se trouve 6 octets plus loin :\\
        \texttt{AD=PEEK(\&H329)*256+PEEK(\&H328)+6}.
  \item Chaîne :\\ \texttt{AD=VARPTR(A\$):AD=PEEK(AD+2)*256+PEEK(AD+1)}\\
        (le descripteur contient la longueur, puis l'adresse, poids faible
        en premier).
  \item Tableau : \texttt{AD=VARPTR(U\%(0))}.
\end{itemize}

\paragraph{Codage dans un tableau d'entiers.} Un entier est rangé poids faible
en premier : les octets \hex{00},\hex{21} s'écrivent donc \hex{2100}
(8448). Il faut souvent ajouter des \texttt{NOP} de cadrage pour que les
données (adresses, longueurs) tombent sur un élément entier. Exemple : une
copie de bloc par \texttt{LDIR} devient
\texttt{NOP / LD HL,source / NOP / LD DE,dest / NOP / LD BC,long / LDIR / RET / NOP} :
\begin{center}\small\ttfamily
\begin{tabular}{@{}lllll@{}}
\toprule
\normalfont Élément & \normalfont Octets & \normalfont Valeur & \normalfont Rôle \\
\midrule
EX(0) & 00 21 & 8448 & NOP ; LD HL,\dots \\
EX(1) & ss ss & -- & adresse source \\
EX(2) & 00 11 & 4352 & NOP ; LD DE,\dots \\
EX(3) & dd dd & -- & adresse destination \\
EX(4) & 00 01 & 256 & NOP ; LD BC,\dots \\
EX(5) & ll ll & -- & longueur \\
EX(6) & ED B0 & -20243 & LDIR \\
EX(7) & C9 00 & 201 & RET ; NOP \\
\bottomrule
\end{tabular}
\end{center}
\nota{Le livre donne $-20234$ pour EX(6) à la p.~19. La bonne valeur,
$-20243$ (\hex{B0ED}), est celle qu'il utilise dans les exemples 4 et 5.}

\subsection{Appel et passage d'arguments \pg{19-21}}

\paragraph{Sans argument.} \texttt{EXEC adresse}. La routine doit se
terminer par \texttt{RET}, sinon le X-07 se bloque et il faut un RESET.

\paragraph{Un argument : \texttt{Z=USR(adresse,argument)}.}
\begin{itemize}
  \item Argument chaîne : \reg{DE} pointe sur le descripteur (longueur, puis
        adresse sur deux octets).
  \item Argument numérique : la valeur est dans l'accumulateur flottant
        (\hex{044E}--\hex{0455}) ; \reg{A} contient le type (2 entier,
        4 simple précision, 8 double précision) et \reg{HL} pointe sur
        \hex{044E}. Pour un entier, il se trouve en \hex{0450}--\hex{0451}.
  \item Valeur de retour : \texttt{USR} renvoie le contenu de
        l'accumulateur.
\end{itemize}

\begin{src}
 & INC HL & HL pointe sur l'entier (\hex{0450}) \\
 & INC HL & \\
 & PUSH HL & sauvegarde pour le retour (inutile si rien à renvoyer) \\
 & LD E,(HL) & \\
 & INC HL & \\
 & LD D,(HL) & \\
 & EX DE,HL & HL = argument entier \\
 & \dots & traitement ; résultat entier dans BC \\
 & POP HL & \\
 & LD (HL),C & renvoi de BC dans l'accumulateur \\
 & INC HL & \\
 & LD (HL),B & \\
 & RET & \\
\end{src}

\paragraph{Plusieurs arguments.} On les range dans un tableau d'entiers dont
on passe l'adresse :
\begin{center}\texttt{T=USR(adresse,VARPTR(ARG(0)))}\end{center}
La routine
récupère cette adresse (comme ci-dessus) puis la place dans \reg{IX} par
\texttt{PUSH DE / POP IX} ; les arguments se lisent alors en
\texttt{(IX+0)}, \texttt{(IX+1)}, etc. Avec une implantation absolue, on peut
aussi déposer les arguments par \texttt{POKE} avant l'appel.

\nota{Le livre écrit que \texttt{ARG(0)=aa00} donne \texttt{(IX+0)=aa} ; en
mémoire, l'octet de poids faible vient en premier : \texttt{(IX+0)=00} et
\texttt{(IX+1)=aa}. Le paragraphe est numéroté 1.1.2.2 au lieu de 1.2.2.2.}

\subsection{Exemples \pg{22-25}}

\subsubsection*{Exemples 1 à 3 : édition automatique de la ligne fautive}
En cas d'erreur, la routine passe directement en mode édition sur la ligne
fautive (\texttt{LIST@}). Le même code de 20 octets est implanté de trois
façons : ligne \texttt{REM}, chaîne de la zone des chaînes, tableau
d'entiers.

\begin{asm}
00 & 06 08    & & LD B,\$08     & \\
02 & E1       & L & POP HL      & vide 8 niveaux de pile \\
03 & 10 FD    & & DJNZ L        & \\
05 & 2A 15 03 & & LD HL,(\$0315) & numéro de la ligne en erreur \\
08 & E5       & & PUSH HL       & \\
09 & D1       & & POP DE        & \\
0A & CD 0B F3 & & CALL \$F30B   & recherche de la ligne \\
0D & FD 21 00 40 & & LD IY,\$4000 & \\
11 & C3 8A FE & & JP \$FE8A     & édition (\texttt{LIST@}) \\
\end{asm}

\begin{verbatim}
10 GOTO 20
15 REM01234567890123456789
20 RESTORE 15:AD%=PEEK(&H329)*256+PEEK(&H328)+6
30 FOR I%=0 TO 19:READ B$:POKE AD%+I%,VAL("&H"+B$):NEXT I%
40 DATA 06,08,E1,10,FD,2A,15,03,E5,D1,CD,0B,F3,FD,21,00,40,C3,8A,FE
50 ON ERROR GOTO 65000
60 ' DEBUT DU PROGRAMME
100 A=5:B=0:C=A/B
64999 END
65000 PRINT "Erreur"
65010 EXEC AD%
\end{verbatim}
Variante « zone des chaînes » (exemple 2) : la ligne 15 disparaît, la ligne 10
devient \texttt{A\$=STRING\$(20,"\#")} et l'adresse est calculée par
\texttt{AD\%=VARPTR(A\$):AD\%=PEEK(AD\%+1)+256*PEEK(AD\%+2)}, en ligne 20
puis de nouveau en ligne 65005, juste avant le \texttt{EXEC}.

Variante « tableau d'entiers » (exemple 3) : mêmes octets, croisés deux à deux.
\begin{verbatim}
10 DIM U%(9)
20 U%(0)=&H0806:U%(1)=&H10E1:U%(2)=&H2AFD:U%(3)=&H0315:U%(4)=&HD1E5
30 U%(5)=&H0BCD:U%(6)=&HFDF3:U%(7)=&H0021:U%(8)=&HC340:U%(9)=&HFE8A
50 ON ERROR GOTO 65000
100 D$=5
64999 END
65000 PRINT"Erreur"
65010 EXEC VARPTR(U%(0))
\end{verbatim}

\subsubsection*{Exemple 4 : remplir un tableau avec une même valeur}
La routine \texttt{LDIR} du §\,1.1 recopie \texttt{A!(0)} dans tout le
tableau : la source et la destination se chevauchent d'un élément. La
longueur vaut \texttt{N*2} pour des entiers, \texttt{N*4} en simple
précision et \texttt{N*8} en double précision.
\begin{verbatim}
10 DEFINT A-Z:N=50:DIM U(7),A!(N)
15 AD=0
20 U(0)=8448:U(2)=4352:U(4)=256:U(6)=-20243:U(7)=201
30 A!(0)=100:GOSUB 100
40 PRINT A!(0);A!(1);A!(20)
50 A!(0)=0:GOSUB 100
60 PRINT A!(0);A!(1);A!(20)
70 END
100 U(1)=VARPTR(A!(0)):U(3)=VARPTR(A!(1)):U(5)=N*4
110 EXEC VARPTR(U(0)):RETURN
\end{verbatim}

\subsubsection*{Exemple 5 : remplir rapidement l'écran vidéo}
Même principe, appliqué à la mémoire vidéo de la X-720 (\hex{8000}). Il faut
ajuster \texttt{US(5)} selon le mode d'écran ; le motif à reproduire est
déposé en \hex{8000} à la ligne 40.
\begin{verbatim}
5 SCREEN 2:CLS
6 LINE (0,0)-(200,180)
10 DEFINT A-Z:DIM US(7)
20 US(0)=8448:US(2)=4352:US(4)=256:US(6)=-20243:US(7)=201
30 US(1)=&H8000:US(3)=&H8001:US(5)=511
40 POKE &H8000,64
50 EXEC VARPTR(US(0))
70 END
\end{verbatim}

\subsubsection*{Exemple 6 : \texttt{RESTORE} calculé}
Le BASIC n'accepte pas de variable après \texttt{RESTORE}. La routine
suivante (18 octets, relogeable) le permet : \texttt{Z=USR(AD,J)} équivaut
à \texttt{RESTORE J}. Elle est utilisée par le logiciel LOGOGÉNÈSE
(§\,\ref{sec:logo}).

\begin{asm}
00 & 23       & & INC HL       & \\
01 & 23       & & INC HL       & \\
02 & 5E       & & LD E,(HL)    & DE = numéro de ligne (argument) \\
03 & 23       & & INC HL       & \\
04 & 56       & & LD D,(HL)    & \\
05 & CD 0D F3 & & CALL \$F30D  & recherche de la ligne : BC = adresse \\
08 & 60       & & LD H,B       & \\
09 & 69       & & LD L,C       & \\
0A & D2 38 F6 & & JP NC,\$F638 & ligne absente : « UL Error » \\
0D & 2B       & & DEC HL       & \\
0E & 22 28 03 & & LD (\$0328),HL & nouveau pointeur des DATA \\
11 & C9       & & RET          & \\
\end{asm}

%----------------------------------------------------------------------
\section{Codage des variables \pg{26-30}}

\subsection{Variables simples}
\texttt{VARPTR(V)} pointe sur la valeur. Les trois octets qui précèdent
contiennent le type et le nom.

\begin{center}\small
\begin{tabular}{@{}>{\raggedright}p{2.6cm}>{\ttfamily\raggedright}p{4.2cm}>{\raggedright\arraybackslash}p{9cm}@{}}
\toprule
Type (taille) & \normalfont Exemple (en hexadécimal) & Codage \\
\midrule
Entier \% (5 octets) & 02 41 00 | D2 04 & \texttt{A\%=1234} : type 2, nom « A », valeur poids faible en premier \\
Simple précision ! (7 octets) & 04 43 44 | 44 12 34 00 & \texttt{CD!=1234} : exposant + \hex{40}, puis 3 octets BCD (0,1234\,E4) \\
Double précision \# (11 octets) & 08 50 49 | 41 31 41 59 26 53 58 98 & \texttt{PI\#=3.1415926535898} : exposant, puis 7 octets BCD \\
Chaîne \$ (6 octets + texte) & 03 54 00 | 04 5B 05 & \texttt{T\$="TOTO"} : longueur, puis adresse du texte \\
\bottomrule
\end{tabular}
\end{center}
Les deux octets de nom contiennent les deux premières lettres ; le second
vaut 0 pour un nom d'une seule lettre.

\subsection{Tableaux}
Devant \texttt{VARPTR(T(0))} se trouve un en-tête : type, nom (2 octets),
taille jusqu'au tableau suivant (2 octets), nombre de dimensions, puis, pour
chaque dimension en partant de la dernière, l'indice maximal + 1 (2 octets).
Les éléments sont rangés en faisant varier le premier indice le plus vite :
\texttt{T(0,0,0)}, \texttt{T(1,0,0)}, \texttt{T(2,0,0)},
\texttt{T(0,1,0)}\dots

\begin{center}\small
\begin{tabular}{@{}l>{\ttfamily}lr@{}}
\toprule
Déclaration & \normalfont En-tête (hexadécimal) & Taille totale \\
\midrule
\texttt{DIM T\%(10)} & 02 54 00 19 00 01 0B 00 & \\
\texttt{DIM T!(10)} & 04 54 00 2F 00 01 0B 00 & \\
\texttt{DIM T\%(4,5)} & 02 54 00 41 00 02 06 00 05 00 & 70 octets \\
\texttt{DIM AT\%(2,4,5)} & 02 41 54 BB 00 03 06 00 05 00 03 00 & 192 octets \\
\texttt{DIM T\%(20)} & & 50 octets \\
\texttt{DIM T!(20)} & & 92 octets \\
\texttt{DIM T\#(20)} & & 176 octets \\
\texttt{DIM T\#(4,5)} & & 250 octets \\
\texttt{DIM T\#(2,4,5)} & & 732 octets \\
\bottomrule
\end{tabular}
\end{center}

%----------------------------------------------------------------------
\section{Routines de la ROM \pg{31-40}}

Ces routines courtes utilisent les adresses des \emph{Mystères du X-07}.
Elles sont données sous forme de source, à implanter où l'on veut.

\subsection{Équivalent de \texttt{INKEY\$} (routine \hex{C90A})}
\begin{src}
INKEY & XOR A & \\
 & CALL \$C90A & lecture d'une touche \\
 & JR Z,INKEY & aucune touche : on recommence \\
 & RET & A = code ASCII \\
\end{src}

\subsection{Tracé de droites (commande \hex{14} du T6834)}
HL pointe sur une table de coordonnées ($X_i$, $Y_i$, $X_f$, $Y_f$) ; B est
le nombre de droites. La routine \hex{C92F} fait avancer HL toute seule.
\begin{src}
DROITE & PUSH BC & \\
 & LD A,\$14 & commande LINE \\
 & LD BC,\$0400 & 4 paramètres, aucune réponse \\
 & CALL \$C92F & \\
 & POP BC & \\
 & DJNZ DROITE & \\
 & RET & \\
\end{src}

\subsection{Caractères graphiques (commande \hex{1A})}
Codes \hex{80}--\hex{9F} et \hex{E0}--\hex{FF}. La table contient, pour chaque
caractère, son code suivi de ses 8 octets de motif.
\begin{src}
CRECAR & LD B,nbre & nombre de caractères \\
 & LD HL,TCAR & table : code, 8 octets, code, \dots \\
CC & PUSH BC & \\
 & LD A,\$1A & commande UDC WRITE \\
 & LD BC,\$0900 & 9 paramètres, aucune réponse \\
 & CALL \$C92F & \\
 & POP BC & \\
 & DJNZ CC & \\
 & RET & \\
\end{src}

\subsection{Touches curseur (commande \hex{02}, STICK)}
\begin{src}
FLECHE & CALL \$C0BD & vide les tampons du clavier \\
 & LD A,\$82 & commande 2 ; bit 7 à 1 : aucun paramètre \\
 & LD BC,\$0001 & 0 paramètre, 1 octet de réponse \\
 & PUSH DE & \\
 & LD DE,BUF & tampon de réponse \\
 & CALL \$C92F & \\
 & LD A,(DE) & \\
 & POP DE & \\
 & CP \$33 & droite \\
 & JR Z,DROITE & \\
 & CP \$37 & gauche \\
 & JR Z,GAUCHE & \\
 & CP \$35 & bas \\
 & JR Z,BAS & \\
 & CP \$31 & haut \\
 & JR Z,HAUT & \\
 & RET & \hex{30} : aucune touche \\
\end{src}
\nota{Les \emph{Mystères} (§\,7.6.1) donnent \hex{32} pour droite et
\hex{36} pour bas ; ce livre utilise \hex{33} et \hex{35}, et il a raison :
le guide technique officiel de la carte XP-140 (tableau~6) donne \hex{31}
haut, \hex{32} haut-droite, \hex{33} droite, \hex{34} bas-droite,
\hex{35} bas, \hex{36} bas-gauche, \hex{37} gauche, \hex{38} haut-gauche.}

\subsection{Test d'une touche (commande \hex{28}, TEST CHR)}
On charge A avec le code de la touche ; au retour, Z\,=\,1 si elle est
enfoncée.
\begin{src}
TOUCHE & LD HL,BUF+4 & \\
 & LD (HL),A & code de la touche \\
 & LD A,\$28 & \\
 & LD BC,\$0101 & 1 paramètre, 1 octet de réponse \\
 & LD DE,BUF+2 & \\
 & CALL \$C92F & \\
 & LD A,(DE) & 0 si enfoncée, sinon \hex{FF} \\
 & OR A & positionne Z \\
 & RET & \\
\end{src}

\subsection{Temporisation}
BC contient la durée.
\begin{src}
DELAY & DEC BC & \\
 & LD A,B & \\
 & OR C & \\
 & JR NZ,DELAY & \\
 & RET & \\
\end{src}
\nota{Le livre écrit \texttt{DEC B} : C n'est alors jamais décrémenté, et la
boucle ne s'arrête jamais si C n'est pas nul. Il faut \texttt{DEC BC}.}

\subsection{Musique (routine \hex{C2DF})}
IX pointe sur la mélodie : note, durée, note, durée\dots{} ; \hex{FF}
termine la mélodie et la note 0 est un silence (bruit). Tous les registres
sont modifiés.
\begin{src}
MUSIC & LD A,(IX+0) & \\
 & CP \$FF & fin de la mélodie ? \\
 & RET Z & \\
 & LD E,A & note \\
 & LD B,(IX+1) & durée \\
 & LD D,0 & \\
 & LD C,D & \\
 & INC IX & \\
 & INC IX & \\
 & LD HL,RTMU & adresse de retour pour C2DF \\
 & PUSH HL & \\
 & PUSH BC & \\
 & PUSH DE & \\
 & POP IY & IY = note (tonalité) \\
 & POP HL & H = durée, L = 0 \\
 & CALL \$C2DF & \\
RTMU & LD HL,(\$0E) & compteur de la note en cours \\
 & LD A,H & \\
 & OR L & \\
 & JR NZ,RTMU & attente de la fin de la note \\
 & JR MUSIC & \\
\end{src}
Mélodie d'exemple :
\begin{verbatim}
AIR1  DB 17,4,17,4,17,4,13,16,0,16,15,4
      DB 15,4,15,4,12,16,0,16,17,4,17,4
      DB 17,4,17,4,13,4,18,4,18,4,18,4
      DB 17,4,25,4,25,4,25,4,22,16,255
\end{verbatim}
\nota{Les commentaires du livre sont inversés (« IY contient la durée »,
« HL la note »). Le code est pourtant conforme aux \emph{Mystères} : C2DF
attend la tonalité dans IY et la durée dans H.}

\subsection{\texttt{PSET}, \texttt{PRESET}, \texttt{POINT} (commandes \hex{11}, \hex{12}, \hex{13})}
\begin{src}
SET & LD A,\$11 & PSET \\
 & JR SET1 & \\
POINT & LD A,\$13 & \\
 & JR SET1 & \\
RESET & LD A,\$12 & PRESET \\
SET1 & LD HL,BUF & BUF contient X et Y \\
 & LD BC,\$0200 & 2 paramètres, aucune réponse \\
 & JP \$C92F & \\
\end{src}
\nota{D'après les \emph{Mystères}, la commande \hex{13} est PEOR (inversion du
point). Le vrai POINT est la commande \hex{10}, qui renvoie un octet : il
faudrait alors \texttt{LD BC,\$0201}.}

\subsection{Équivalent de \texttt{INPUT "message";A\$}}
La routine \hex{FEF7} affiche le message, \hex{EBF2} lit le clavier. La
chaîne saisie (20 caractères au plus) est copiée dans BUF et sa longueur
dans LONG ; une réponse vide fait recommencer.
\begin{src}
ACQMES & LD HL,MESSAGE & message terminé par 0 \\
 & CALL \$FEF7 & \\
 & CALL \$EBF2 & saisie ; HL = tampon $-1$ \\
 & INC HL & \\
 & LD DE,BUF & \\
 & LD C,0 & \\
ACQ & LD A,(HL) & \\
 & LD (DE),A & \\
 & OR A & \\
 & JR Z,ACQ2 & fin de la saisie \\
 & INC HL & \\
 & INC DE & \\
 & INC C & \\
 & LD A,C & \\
 & CP 20 & \\
 & JR C,ACQ & \\
ACQ2 & LD A,C & \\
 & CP 0 & \\
 & JR Z,ACQMES & chaîne vide : on recommence \\
 & LD (LONG),A & \\
 & RET & \\
\end{src}

\subsection{Équivalent de \texttt{INPUT "message";A\%}}
\begin{src}
ACQNUM & LD HL,MESSAGE & \\
 & CALL \$FEF7 & \\
 & CALL \$EBF2 & saisie \\
 & RST \$10 & premier caractère \\
 & CALL \$F595 & conversion ASCII $\to$ entier dans DE \\
\end{src}
\nota{Vérifié dans l'émulateur (vraie ROM) : \hex{F595} est la routine qui lit
un \textbf{numéro de ligne}. Elle ne lit qu'un entier \textbf{non signé} de
0 à 65535 : une saisie de \texttt{-56} donne 0, et \texttt{40000} est accepté
sans erreur \texttt{OV}. Ce n'est donc un équivalent de
\texttt{INPUT A\%} que pour les nombres positifs. Les espaces de tête sont
bien ignorés (RST 10).}

\subsection{Sortie sur la table traçante X-710 (routine \hex{CEF7})}
\begin{src}
 & LD HL,CHAINE & \\
LPRINT & LD A,(HL) & \\
 & OR A & fin de chaîne ? \\
 & RET Z & \\
 & PUSH HL & \\
 & CALL \$CEF7 & envoi à la X-710 \\
 & POP HL & \\
 & INC HL & \\
 & JR LPRINT & \\
\end{src}
Chaîne d'exemple, en décimal :
\texttt{18,83,48,13,67,49,77,48,45,54,13,65,13,10,00}. Le code 18 fait
passer en mode graphique ; « S0 » règle la taille, « C1 » la couleur ;
« M0,-6 » remonte de 6 points ; 65 est la lettre A ; 13 et 10 sont CR et LF ;
0 termine la chaîne.

%----------------------------------------------------------------------
\section{Extension des fonctions du BASIC \pg{41-45}}

\subsection{Le crochet d'erreur}
Avant d'afficher une erreur, la ROM exécute l'instruction placée en
\hex{00AB} (normalement \texttt{RET}, \hex{C9}). On y met un
\texttt{JP xxxx} vers sa propre routine, qui reçoit le numéro d'erreur dans
E. Si elle ne traite pas l'erreur, elle rend la main par \texttt{RET} et
l'erreur suit son cours normal. Codes utiles : 2 = erreur de syntaxe,
6 = débordement, 8 = numéro de ligne indéfini.

\subsection{Messages d'erreur en français}
\fiche{\hex{1C00}--\hex{1C8C}}{141 octets}{\texttt{EXEC \&H1C00} (installe le détournement)}

\begin{asm}
1C00 & 3E C3    & & LD A,\$C3      & code de JP \\
1C02 & 32 AB 00 & & LD (\$00AB),A  & \\
1C05 & 21 0C 1C & & LD HL,DB       & \\
1C08 & 22 AC 00 & & LD (\$00AC),HL & JP DB en \hex{00AB} \\
1C0B & C9       & & RET            & \\
1C0C & 7B       & DB & LD A,E      & numéro d'erreur \\
1C0D & 32 00 1B & & LD (\$1B00),A  & (absent du source du livre) \\
1C10 & FE 02    & & CP 2           & syntaxe ? \\
1C12 & 28 09    & & JR Z,E2        & \\
1C14 & FE 06    & & CP 6           & débordement ? \\
1C16 & 28 11    & & JR Z,E6        & \\
1C18 & FE 08    & & CP 8           & ligne indéfinie ? \\
1C1A & 28 19    & & JR Z,E8        & \\
1C1C & C9       & & RET            & autre erreur : traitement normal \\
1C1D & 21 41 1C & E2 & LD HL,SY    & \\
1C20 & CD F7 FE & & CALL \$FEF7    & affiche le message \\
1C23 & CD 79 1C & & CALL NO        & affiche « LIGNE n » \\
1C26 & C3 3D F2 & & JP \$F23D      & arrêt, attente curseur \\
1C29 & \dots    & E6 & \dots       & idem avec DE (\hex{1C56}) \\
1C35 & \dots    & E8 & \dots       & idem avec ID (\hex{1C67}) \\
1C41 & \dots    & SY & DEFM        & « ERREUR DE SYNTAXE », CR LF, 0 \\
1C56 & \dots    & DE & DEFM        & « DEBORDEMENT EN », CR LF, 0 \\
1C67 & \dots    & ID & DEFM        & « LIGNE INDEFINIE », CR LF, 0 \\
1C79 & 21 86 1C & NO & LD HL,N1    & \\
1C7C & CD F7 FE & & CALL \$FEF7    & « LIGNE » \\
1C7F & 2A DB 01 & & LD HL,(\$01DB) & dernière ligne exécutée \\
1C82 & CD 98 BB & & CALL \$BB98    & affichage décimal \\
1C85 & C9       & & RET            & \\
1C86 & \dots    & N1 & DEFM        & « LIGNE », 0 \\
\end{asm}

\begin{verbatim}
10 REM Messages d'erreurs en francais
20 FOR I=&H1C00 TO &H1C00+140
30 READ A$:POKE I,VAL("&H"+A$)
40 NEXT
50 PRINT "EXEC &H1C00 sera necessaire au detournement des futures erreurs"
60 DATA 3E,C3,32,AB,0,21,C,1C,22,AC,0,C9,7B,32,0,1B,FE,2,28,9,FE,6,28,11,FE,8
70 DATA 28,19,C9,21,41,1C,CD,F7,FE,CD,79,1C,C3,3D,F2,21,56,1C,CD,F7,FE,CD,79,1C
80 DATA C3,3D,F2,21,67,1C,CD,F7,FE,CD,79,1C,C3,3D,F2,45,52,52,45,55,52,20,44
90 DATA 45,20,53,59,4E,54,41,58,45,20,A,D,0,44,45,42,4F,52,44,45,4D,45,4E,54,20
100 DATA 45,4E,A,D,0,4C,49,47,4E,45,20,49,4E,44,45,46,49,4E,49,45,A,D,0,21,86
110 DATA 1C,CD,F7,FE,2A,DB,1,CD,98,BB,C9,4C,49,47,4E,45,20,0
\end{verbatim}
\nota{Les DATA du livre contiennent trois fautes dans les messages :
« ERREURE DE SYMTAXE », et un code \hex{04} au lieu de \hex{44} (« D ») dans
« DEBORDEMENT ». Elles sont corrigées ci-dessus sans changer la longueur
des messages (« ERREUR DE SYNTAXE » suivi d'une espace), donc sans décaler
les adresses.}

\subsection{Nouvelle instruction \texttt{CARTE}}
Taper \texttt{CARTE} produit une erreur de syntaxe (E\,=\,2). La routine
relit alors le mot à rebours (HL pointe juste après) et, s'il s'agit bien de
\texttt{CARTE}, saute en \hex{C4D1} (« CREATE SYSTEM ? », effacement de la
carte mémoire).
\fiche{\hex{1C00}--\hex{1C33}}{52 octets}{\texttt{EXEC \&H1C00}, puis taper \texttt{CARTE}}
\begin{asm}
1C00 & 3E C3 \dots & & \dots         & installation du JP en \hex{00AB} (comme ci-dessus) \\
1C0C & 7B       & DB & LD A,E        & \\
1C0D & FE 02    & & CP 2             & erreur de syntaxe ? \\
1C0F & C0       & & RET NZ           & \\
1C10 & CD 31 1C & & CALL PR          & caractère précédent \\
1C13 & FE 45    & & CP "E"           & \\
1C15 & C0       & & RET NZ           & \\
 & \dots & & \dots & idem pour T, R, A, C \\
1C2E & C3 D1 C4 & & JP \$C4D1       & « CREATE SYSTEM ? » \\
1C31 & 2B       & PR & DEC HL        & \\
1C32 & 7E       & & LD A,(HL)        & \\
1C33 & C9       & & RET              & \\
\end{asm}
\begin{verbatim}
10 ' Instruction CARTE
20 FOR I=&H1C00 TO &H1C00+51
30 READ A$:POKE I,VAL("&H"+A$)
40 NEXT
50 EXEC &H1C00
60 DATA 3E,C3,32,AB,0,21,C,1C,22,AC,0,C9,7B,FE,2,C0,CD,31,1C,FE,45,C0,CD,31,
   1C,FE,54,C0,CD,31,1C,FE,52,C0,CD,31,1C,FE,41,C0,CD,31,1C,FE,43,C0,C3,D1,
   C4,2B,7E,C9
\end{verbatim}

\subsection{Détournement de \texttt{PAINT}}
Sans la X-720, \texttt{PAINT} passe par le saut placé en \hex{0099} (qui
mène à SN ERROR). On peut le détourner de la même manière :
\begin{src}
 & ORG \$1C00 & \\
 & LD A,\$C3 & \\
 & LD (\$99),A & \\
 & LD HL,DB & \\
 & LD (\$9A),HL & JP DB en \hex{0099} \\
 & RET & \\
DB & PUSH HL & \\
 & LD HL,ME & \\
 & CALL \$D5B0 & affiche « PAINT » \\
 & POP HL & \\
 & RET & \\
ME & DEFM PAINT & \\
 & DEFB \$00 & \\
\end{src}

%======================================================================
\part*{2\textsuperscript{e} partie : Hard \pg{46-82}}
\addcontentsline{toc}{part}{2e partie : Hard}

%----------------------------------------------------------------------
\section{Interruptions et entrées/sorties \pg{47-68}}

\subsection{Les interruptions du X-07 \pg{47-51}}

Une interruption permet à un périphérique de signaler un événement sans que
le programme ait à le surveiller en permanence. Le Z80 la prend en compte à
la fin de l'instruction en cours : il empile l'adresse de retour (PC),
exécute la routine de traitement puis reprend. \texttt{DI} et \texttt{EI}
interdisent et autorisent les interruptions masquables.

\begin{center}\small
\begin{tabular}{@{}l>{\raggedright}p{5.5cm}>{\ttfamily}l>{\ttfamily}l@{}}
\toprule
Ligne & Source & \normalfont Vecteur & \normalfont Routine ROM \\
\midrule
INT A & sous-processeur T6834 (clavier\dots) & 003C & JP C799 \\
INT B & interface série & 0034 & JP C7A3 \\
INT C & X-720 seulement : horloge du curseur & 002C & RET (C9) \\
NMI & non utilisée sur le X-07 ; non masquable, priorité maximale & 0066 & RET (C9) \\
\bottomrule
\end{tabular}
\end{center}
Ces vecteurs sont en RAM (zone système) : on peut les modifier. NMI, qui ne
fait rien, est la plus facile à détourner.

\nota{Au §\,5.1.4, le livre dit que la ligne B vient du sous-processeur, ce
qui contredit le §\,5.1.2 (INT A = T6834, INT B = série).}

\paragraph{Déclencher une NMI à la main \pg{51-53}.} Les bornes INT A et NMI
sont accessibles sur le connecteur d'extension arrière (brochage p.~127 du
guide utilisateur). NMI (broche 20) est active à l'état haut, INT A à l'état
bas. Le +5\,V est sur les broches 39 et 40. Un bref contact entre les
broches 20 et 39 déclenche une NMI. Les auteurs conseillent un bouton-poussoir
entre ces deux broches : un contact trop long fausse le fonctionnement.

\begin{center}\small\ttfamily
\begin{tabular}{@{}l*{10}{c}@{}}
\toprule
\normalfont Rangée haute (paires) & 2 & 4 & \dots & \textbf{20} NMI & \dots & & & & & 40 +5V \\
\normalfont Rangée basse (impaires) & 1 & 3 & \dots & 19 & \dots & & & & & \textbf{39} +5V \\
\bottomrule
\end{tabular}
\end{center}

\paragraph{Pour aller plus loin.} En divisant le signal CLOCK (broche 29)
jusqu'à 5 à 10\,Hz et en le renvoyant sur la broche 20, on obtient un
« top » périodique qui permet un multitâche simple. Un montage anti-rebond
(bascule) est recommandé. Applications envisagées : détection de niveaux,
surveillance, sonnerie de téléphone, disquettes, DMA.

\subsubsection*{Expérience 1 : un BEEP à chaque NMI}
Le BASIC compte de 1 à 10\,000 pendant que chaque NMI émet un bip
(caractère BELL, code 7).
\fiche{\hex{1B00}--\hex{1B12}}{19 octets}{\texttt{EXEC \&H1B00}, puis relier brièvement les broches 20 et 39}
\begin{asm}
1B00 & 3E C3    & & LD A,\$C3      & \\
1B02 & 32 66 00 & & LD (\$0066),A  & \\
1B05 & 21 0C 1B & & LD HL,DB       & \\
1B08 & 22 67 00 & & LD (\$0067),HL & JP DB en \hex{0066} \\
1B0B & C9       & & RET            & \\
1B0C & 3E 07    & DB & LD A,\$07   & BELL \\
1B0E & CD BE C1 & & CALL \$C1BE    & sortie console \\
1B11 & ED 4D    & & RETI           & \\
\end{asm}
\begin{verbatim}
5 REM LE BEEP
10 FOR I=&H1B00 TO &H1B12
20 READ A$:POKE I,VAL("&H"+A$)
30 NEXT
35 EXEC &H1B00
40 CLS:PRINT"COURT-CIRCUITEZ LES BROCHES"
50 FOR I=1TO10000
60 LOCATE10,2:PRINTI
70 NEXT
100 DATA 3E,C3,32,66,00,21,C,1B
101 DATA 22,67,0,C9
102 DATA 3E,7,CD,BE,C1,ED,4D
\end{verbatim}
\nota{Le listing du livre ne contient pas de \texttt{EXEC \&H1B00} : la ligne
35 est ajoutée ici. Une NMI doit normalement se terminer par \texttt{RETN}
(\hex{ED 45}), qui restaure l'état des interruptions ; le livre emploie
\texttt{RETI} (\hex{ED 4D}).}

\nota{Vérifié dans l'émulateur (vraie ROM) : la routine écrase le registre A
(\texttt{LD A,\$07}) sans le sauvegarder. Le BASIC interrompu est faussé :
avec des NMI répétées, le compteur repart en arrière et un \texttt{Break in
60} intempestif arrête le programme. Il faut encadrer la routine par
\texttt{PUSH AF} et \texttt{POP AF} :
\texttt{1B0C F5 / 3E 07 / CD BE C1 / F1 / ED 45}, soit
\texttt{102 DATA F5,3E,7,CD,BE,C1,F1,ED,45} et \texttt{FOR I=\&H1B00 TO
\&H1B14} en ligne 10. Ainsi corrigée, chaque NMI donne un bip et le comptage
reste juste. L'expérience 2, qui sauvegarde tous ses registres, fonctionne
telle quelle.}

\subsubsection*{Expérience 2 : deux routines en langage machine}
Une boucle affiche les codes 0 à 255 ; chaque NMI affiche
« INTERUPTION !!!! » sans arrêter la boucle.
\fiche{\hex{1A00}--\hex{1A0B} et \hex{1B00}--\hex{1B4C}}{12 + 77 octets}{\texttt{EXEC \&H1A00:EXEC \&H1B00}}
\begin{asm}
1A00 & 3E C3    & & LD A,\$C3      & installation : \\
1A02 & 32 66 00 & & LD (\$0066),A  & JP \hex{1B20} en \hex{0066} \\
1A05 & 21 20 1B & & LD HL,\$1B20   & \\
1A08 & 22 67 00 & & LD (\$0067),HL & \\
1A0B & C9       & & RET            & \\
1B00 & 3E 00    & & LD A,0         & 1\textsuperscript{er} programme \\
1B02 & 3C       & B1 & INC A       & \\
1B03 & 21 02 05 & & LD HL,\$0502   & position du curseur \\
1B06 & 22 B8 00 & & LD (\$00B8),HL & \\
1B09 & CD BE C1 & & CALL \$C1BE    & affiche le caractère A \\
1B0C & CD 11 1B & & CALL TP        & \\
1B0F & 18 F1    & & JR B1          & \\
1B11 & F5       & TP & PUSH AF     & temporisation \\
1B12 & C5       & & PUSH BC        & \\
1B13 & 3E 7F    & & LD A,\$7F      & \\
1B15 & 06 FF    & B2 & LD B,\$FF   & \\
1B17 & 05       & B3 & DEC B       & \\
1B18 & 20 FD    & & JR NZ,B3       & \\
1B1A & 3D       & & DEC A          & \\
1B1B & 20 F8    & & JR NZ,B2       & \\
1B1D & C1       & & POP BC         & \\
1B1E & F1       & & POP AF         & \\
1B1F & C9       & & RET            & \\
1B20 & F5 C5 D5 E5 & DB & PUSH AF\dots HL & 2\textsuperscript{e} programme (NMI) \\
1B24 & 21 01 01 & & LD HL,\$0101   & \\
1B27 & 22 B8 00 & & LD (\$00B8),HL & \\
1B2A & 21 3C 1B & & LD HL,ME       & \\
1B2D & CD F7 FE & & CALL \$FEF7    & affiche le message \\
1B30 & CD 11 1B & & CALL TP        & \\
1B33 & CD 9E CE & & CALL \$CE9E    & efface l'écran \\
1B36 & E1 D1 C1 F1 & & POP HL\dots AF & \\
1B3A & ED 4D    & & RETI           & \\
1B3C & \dots    & ME & DEFM        & « INTERUPTION !!!! », 0 \\
\end{asm}
\nota{Vérifié dans la ROM : \hex{00B8} contient la ligne et \hex{00B9} la
colonne, comptées à partir de 1. \texttt{LD HL,\$0502 : LD (\$00B8),HL}
place donc le curseur en colonne 5, ligne 2, et \texttt{\$0101} le ramène en
haut à gauche.}
\begin{verbatim}
10 REM LES 2
20 FOR I=&H1A00 TO &H1A0B
30 READ A$:POKE I,VAL("&H"+A$):NEXT
50 RESTORE 700
60 FOR I=&H1B00 TO &H1B4C
70 READ A$:POKE I,VAL("&H"+A$):NEXT
90 EXEC&H1A00:EXEC&H1B00
500 DATA 3E,C3,32,66,00,21,20,1B,22
501 DATA 67,0,C9
700 DATA 3E,0,3C,21,2,5,22,B8,0
701 DATA CD,BE,C1,CD,11,1B
710 DATA 18,F1,F5,C5,3E,7F,6,FF,5
711 DATA 20,FD,3D,20,F8
720 DATA C1,F1,C9,F5,C5,D5,E5,21,1
721 DATA 1,22,B8,0,21
730 DATA 3C,1B,CD,F7,FE,CD,11,1B,CD
731 DATA 9E,CE,E1,D1
740 DATA C1,F1,ED,4D,49,4E,54,45,52
741 DATA 55,50,54,49,4F,4E,20,21
750 DATA 21,21,21,0
\end{verbatim}

%----------------------------------------------------------------------
\subsection{Les entrées/sorties \pg{57-68}}

Le X-07 a quatre prises : magnétophone, série, parallèle et extension.
Plutôt que de programmer directement les circuits, le livre réutilise les
routines de la ROM qui servent à \texttt{INIT \#n} (ouverture d'un fichier
d'entrée/sortie).

\paragraph{Les dispositifs.}
\begin{center}\small
\begin{tabular}{@{}l>{\ttfamily}l>{\raggedright\arraybackslash}p{10cm}@{}}
\toprule
Nom & \normalfont Table & Rôle et initialisation \\
\midrule
KBD & E778 & clavier et écran LCD ; l'ouverture vide le tampon du clavier \\
CASI & E787 & entrée cassette \\
CASO & E796 & sortie cassette \\
COM & E7A4 & prise série ; vitesse et format à fournir \\
GPR & E7B2 & table traçante X-710 (taille 2 ; CR et LF envoyés) \\
PRT & E7C0 & imprimante série \\
RAM & E7CE & fichier en RAM (créé s'il n'existe pas, avec la taille demandée) \\
OPT & E7DC & coupleur optique \\
CON & E7EA & console (l'ouverture ne fait rien) \\
LPT & E7F8 & imprimante parallèle Centronics (CR et LF) \\
\bottomrule
\end{tabular}
\end{center}

Chaque entrée de la table de la ROM contient : le nom (par exemple
« CASI: »), le numéro du dispositif, puis les adresses des sous-programmes
de sortie d'un octet, d'entrée d'un octet, d'initialisation, et
éventuellement un sous-programme optionnel. L'adresse indiquée ci-dessus
pointe sur le numéro du dispositif. Exemple pour CASI :
\texttt{43 41 53 49 BA | 00 02 | 12 E5 | 0A E0 | 35 DE | 0F E1}.

\paragraph{Routines utilisées.}
\begin{center}\small
\begin{tabular}{@{}>{\ttfamily}l>{\raggedright\arraybackslash}p{12cm}@{}}
\toprule
\normalfont Adresse & Rôle \\
\midrule
E6A8 & ouverture : DE = table du dispositif, IY = entrée de la table des fichiers ouverts (\hex{02C5} pour le 1\textsuperscript{er}, \hex{02CD} pour le 2\textsuperscript{e}), A = 0. Pour COM, B = mode, IX = vitesse \\
E827 & assignation : DE = \hex{02C5} (ou \hex{02CD}) choisit le fichier courant \\
E88F & sortie de A vers le fichier courant (= RST \hex{28}) \\
E8D4 & entrée d'un octet du fichier courant dans A (ne rend la main que si une donnée est présente ; contient l'appel ABORT) \\
\bottomrule
\end{tabular}
\end{center}
La table des fichiers ouverts (\hex{02C5}--\hex{02E8}) contient 5 entrées de
8 octets : adresse du dispositif et sauvegarde de HL, DE et BC. On peut donc
ouvrir 5 fichiers au plus.

\subsubsection*{Sortie sur cassette (CASO)}
Le moteur du magnétophone n'est pas commandé automatiquement : il faut mettre
à 1 le bit 0 (REM) du port \hex{F4}. Les données sont précédées de trois
\hex{00} et suivies de trois \hex{FF}.
\fiche{\hex{1C00}--\hex{1C67}}{104 octets (données de \hex{1C61} à \hex{1C67})}{\texttt{EXEC \&H1C00}}
\begin{asm}
1C00 & 11 96 E7 & & LD DE,\$E796   & table CASO \\
1C03 & FD 21 C5 02 & & LD IY,\$02C5 & fichier \#1 \\
1C07 & 3E 00    & & LD A,0         & \\
1C09 & CD A8 E6 & & CALL \$E6A8    & INIT \#1,"CASO:" \\
1C0C & DB F4    & & IN A,(\$F4)    & \\
1C0E & CB C7    & & SET 0,A        & moteur en marche \\
1C10 & D3 F4    & & OUT (\$F4),A   & \\
1C12 & CD 54 1C & & CALL TP        & temporisation \\
1C15 & CD 54 1C & & CALL TP        & \\
1C18 & 11 C5 02 & & LD DE,\$02C5   & \\
1C1B & CD 27 E8 & & CALL \$E827    & assignation \\
1C1E & \dots    & & 3 $\times$ (LD A,0 ; CALL \$E88F) & trois \hex{00} de début \\
1C2D & 21 61 1C & & LD HL,DD       & données \\
1C30 & 7E       & ST & LD A,(HL)   & \\
1C31 & FE 00    & & CP 0           & fin des données ? \\
1C33 & 28 06    & & JR Z,TE        & \\
1C35 & CD 8F E8 & & CALL \$E88F    & envoi de l'octet \\
1C38 & 23       & & INC HL         & \\
1C39 & 18 F5    & & JR ST          & \\
1C3B & \dots    & TE & 3 $\times$ (LD A,\$FF ; CALL \$E88F) & trois \hex{FF} de fin \\
1C4A & CD 54 1C & & CALL TP        & \\
1C4D & DB F4    & & IN A,(\$F4)    & \\
1C4F & CB 87    & & RES 0,A        & moteur arrêté \\
1C51 & D3 F4    & & OUT (\$F4),A   & \\
1C53 & C9       & & RET            & \\
1C54 & C5       & TP & PUSH BC     & temporisation \\
1C55 & 06 7F    & & LD B,\$7F      & \\
1C57 & 0E FF    & E1 & LD C,\$FF   & \\
1C59 & 0D       & E2 & DEC C       & \\
1C5A & 20 FD    & & JR NZ,E2       & \\
1C5C & 05       & & DEC B          & \\
1C5D & 20 F8    & & JR NZ,E1       & \\
1C5F & C1       & & POP BC         & \\
1C60 & C9       & & RET            & \\
1C61 & 01 \dots 06 00 & DD & DEFB    & données 1 à 6, puis 0 \\
\end{asm}
\begin{verbatim}
10 FOR I=&H1C00 TO &H1C00+103
20 READ A$:POKE I,VAL("&H"+A$)
30 NEXT I
40 PRINT "EXEC &H1C00 POUR DEMARRER LE PROGRAMME"
50 DATA 11,96,E7,FD,21,C5,2,3E,0,CD,A8,E6,DB,F4,CB,C7,D3,F4,CD,54,1C,CD,54,1C,
   11,C5,2,CD,27,E8,3E,0,CD,8F,E8,3E,0,CD,8F,E8,3E,0,CD,8F,E8,21,61,1C,7E,
   FE,0,28,6,CD,8F,E8,23,18,F5,3E,FF,CD,8F,E8,3E,FF,CD,8F,E8,3E,FF,CD,8F,E8,
   CD,54,1C,DB,F4,CB,87,D3,F4,C9,C5,6,7F,E,FF,D,20,FD,5,20,F8,C1,C9,1,2,3,4,
   5,6,0
\end{verbatim}
\nota{Le premier listing de la p.~60 donne \texttt{LD DE,\$E987} pour la
table CASO : c'est \hex{E796}. Le source porte \texttt{LD B,\$00} dans la
temporisation, les DATA \texttt{LD B,\$7F}. Comme \hex{00} sert de marqueur
de fin, les données ne peuvent pas contenir d'octet nul.}

\subsubsection*{Lecture sur cassette (CASI)}
La routine attend les trois \hex{00} de début, puis range les octets à partir
de \hex{1B00} jusqu'aux trois \hex{FF} de fin. Aucune erreur de lecture n'est
traitée.
\begin{src}
 & LD DE,\$E787 & INIT \#1,"CASI:" \\
 & LD IY,\$02C5 & \\
 & LD A,0 & \\
 & CALL \$E6A8 & \\
 & IN A,(\$F4) & moteur en marche \\
 & SET 0,A & \\
 & OUT (\$F4),A & \\
 & LD DE,\$02C5 & assignation \\
 & CALL \$E827 & \\
E1 & CALL \$E8D4 & attente des trois \hex{00} \\
 & JR NZ,E1 & (répété 3 fois) \\
 & LD HL,\$1B00 & zone de réception \\
E2 & CALL \$E8D4 & \\
 & CP \$FF & \\
 & JR Z,FI & \\
 & LD (HL),A & \\
 & INC HL & \\
 & JR E2 & \\
FI & CALL \$E8D4 & 2\textsuperscript{e} \hex{FF} ? \\
 & \dots & (idem pour le 3\textsuperscript{e}) \\
F2 & IN A,(\$F4) & moteur arrêté \\
 & RES 0,A & \\
 & OUT (\$F4),A & \\
 & RET & \\
\end{src}
\begin{verbatim}
10 FOR I=&H1C00 TO &H1C00+81
20 READ A$:POKE I,VAL("&H"+A$)
30 NEXT
40 PRINT "EXEC &H1C00 pour demarrer ..."
50 DATA 11,87,E7,FD,21,C5,2,3E,0,CD,A8,E6,DB,F4,CB,C7,D3,F4,11,C5,2,CD,27,E8,
   CD,D4,E8,20,FB,CD,D4,E8,20,F6,CD,D4,E8,20,F1,21,0,1B,CD,D4,E8,FE,FF,28,4,
   77,23,18,F5,CD,D4,E8,FE,FF,28,4,7E,23,18,EA,CD,D4,E8,FE,FF,28,4,7E,23,18,
   DF,DB,F4,CB,87,D3,F4,C9
\end{verbatim}
\nota{Les DATA du livre contiennent un octet parasite \hex{8F} après le
troisième \texttt{CALL \$E8D4} : le dernier \texttt{JR} saute alors au milieu
d'une instruction et la machine plante. Il est retiré ci-dessus (82 octets,
\texttt{FOR \dots +81}). La routine reste fragile : un \hex{FF} isolé dans
les données est remplacé par l'octet en mémoire (\texttt{LD A,(HL)} au lieu
d'un rangement).}

\subsubsection*{Clavier vers imprimante (KBD et GPR/LPT)}
Chaque touche frappée est envoyée à l'imprimante ; CTRL N arrête.
\fiche{\hex{1C00}--\hex{1C30}}{49 octets}{\texttt{EXEC \&H1C00}}
\begin{asm}
1C00 & 11 F8 E7 & & LD DE,\$E7F8   & imprimante (voir N.d.R.) \\
1C03 & FD 21 C5 02 & & LD IY,\$02C5 & fichier \#1 \\
1C07 & 3E 00    & & LD A,0         & \\
1C09 & CD A8 E6 & & CALL \$E6A8    & \\
1C0C & 11 78 E7 & & LD DE,\$E778   & KBD \\
1C0F & FD 21 CD 02 & & LD IY,\$02CD & fichier \#2 \\
1C13 & 3E 00    & & LD A,0         & \\
1C15 & CD A8 E6 & & CALL \$E6A8    & \\
1C18 & 11 CD 02 & BL & LD DE,\$02CD & clavier \\
1C1B & CD 27 E8 & & CALL \$E827    & \\
1C1E & CD D4 E8 & & CALL \$E8D4    & attente d'une touche \\
1C21 & F5       & & PUSH AF        & \\
1C22 & 11 C5 02 & & LD DE,\$02C5   & imprimante \\
1C25 & CD 27 E8 & & CALL \$E827    & \\
1C28 & F1       & & POP AF         & \\
1C29 & FE 0E    & & CP \$0E        & CTRL N : fin \\
1C2B & C8       & & RET Z          & \\
1C2C & CD 8F E8 & & CALL \$E88F    & impression \\
1C2F & 18 E7    & & JR BL          & \\
\end{asm}
\begin{verbatim}
10 FOR I=&H1C00 TO &H1C00+48
20 READ A$:POKE I,VAL("&H"+A$):NEXT
40 PRINT "EXEC &H1C00 pour demarrer ..."
50 DATA 11,F8,E7,FD,21,C5,2,3E,0,CD,A8,E6,11,78,E7,FD,21,CD,2,3E,0,CD,A8,E6,
   11,CD,2,CD,27,E8,CD,D4,E8,F5,11,C5,2,CD,27,E8,F1,FE,E,C8,CD,8F,E8,18,E7
\end{verbatim}
\nota{Le source vise la X-710 (GPR, \hex{E7B2}) et appelle par erreur
\texttt{\$F6A8} ; les DATA utilisent \hex{E7F8} (LPT, imprimante parallèle)
et \hex{E6A8}. Pour la X-710, remplacer \texttt{F8,E7} par \texttt{B2,E7}.}

\subsubsection*{Prise série : émission vers un Minitel}
Pour COM, l'ouverture prend aussi le mode de l'ACIA dans B (ici « G ») et la
vitesse dans IX (1200 bauds = \hex{04B0}). Le code 12 (\hex{0C}) efface
l'écran du Minitel.
\fiche{\hex{1C00}--\hex{1C3F}}{64 octets}{\texttt{EXEC \&H1C00}}
\begin{asm}
1C00 & 11 A4 E7 & & LD DE,\$E7A4   & table COM \\
1C03 & 3E 00    & & LD A,0         & \\
1C05 & 06 47    & & LD B,"G"       & mode G de l'ACIA \\
1C07 & DD 21 B0 04 & & LD IX,1200  & vitesse \\
1C0B & FD 21 C5 02 & & LD IY,\$02C5 & \\
1C0F & CD A8 E6 & & CALL \$E6A8    & \\
1C12 & 11 C5 02 & & LD DE,\$02C5   & \\
1C15 & CD 27 E8 & & CALL \$E827    & assignation \\
1C18 & 21 27 1C & & LD HL,ME       & \\
1C1B & 7E       & ET & LD A,(HL)   & \\
1C1C & FE 00    & & CP 0           & \\
1C1E & 28 06    & & JR Z,FI        & \\
1C20 & CD 8F E8 & & CALL \$E88F    & émission \\
1C23 & 23       & & INC HL         & \\
1C24 & 18 F5    & & JR ET          & \\
1C26 & C9       & FI & RET         & \\
1C27 & 0C \dots & ME & DEFB 12 ; DEFM & effacement, « CA MARCHE BIEN COMME CA », 0 \\
\end{asm}
\begin{verbatim}
20 FOR I=&H1C00 TO &H1C00+63
30 READ A$:POKE I,VAL("&H"+A$):NEXT
50 PRINT " EXEC &H1C00 pour demarrer ..."
60 DATA 11,A4,E7,3E,0,06,47,DD,21,B0,4,FD,21,C5,2,CD,A8,E6,11,C5,2,CD,27,E8,
   21,27,1C,7E,FE,0,28,6,CD,8F,E8,23,18,F5,C9,C,43,41,20,4D,41,52,43,48,45,
   20,42,49,45,4E,20,43,4F,4D,4D,45,20,43,41,0
\end{verbatim}

\subsubsection*{Prise série : réception}
Même ouverture, puis lecture et affichage d'un caractère.
\begin{verbatim}
20 FOR I=&H1C00 TO &H1C00+30
30 READ A$:POKE I,VAL("&H"+A$):NEXT
50 PRINT " EXEC &H1C00 pour demarrer ..."
60 DATA 11,A4,E7,3E,0,6,47,DD,21,B0,4,FD,21,C5,2,CD,A8,E6,11,C5,2,CD,27,E8,
   CD,D4,E8,CD,BE,C1,C9
\end{verbatim}
\nota{Le livre écrit \texttt{FOR \dots +29} : le \texttt{RET} final
(31\textsuperscript{e} octet) n'est pas chargé. Il écrit aussi
\texttt{CD,DE,C1} au lieu de \texttt{CD,BE,C1} (\texttt{CALL \$C1BE}). Les
deux erreurs sont corrigées ci-dessus. La routine ne lit qu'un seul
caractère ; pour une réception continue, il faut boucler sur
\texttt{CALL \$E8D4 / CALL \$C1BE}.}

%----------------------------------------------------------------------
\section{L'interface vidéo X-720 \pg{69-82}}

\subsection{Matériel}
La X-720 relie le X-07 à un téléviseur par la prise Péritel ou par un
modulateur. Les auteurs la jugent encombrante, non autonome, lente et peu
servie en logiciels.

\begin{itemize}
  \item Générateur vidéo \textbf{Motorola MC6847}.
  \item 6\,Ko de RAM vidéo (\hex{8000}--\hex{97FF}), soit trois circuits
        6117 de 2\,Ko (IC203, IC204, IC205). Deux mémorisent l'état des
        points, le troisième les attributs (couleur).
  \item ROM : générateur de caractères et 4\,Ko de complément du BASIC
        (\hex{A000}--\hex{AFFF}).
  \item Oscillateur couleur à 3,579545\,MHz ; alimentation par adaptateur
        AD-5 et batterie NiCd.
  \item Connecteur vers le X-07 : NMI broche 20, CLK broche 29, VCC2
        broches 39--40 (confirmé par le schéma de la fig.~11).
\end{itemize}

\paragraph{Ports d'entrée/sortie (fig.~9).}
\begin{center}\small
\begin{tabular}{@{}>{\ttfamily}l l *{8}{c}@{}}
\toprule
\normalfont Ports & Sens & D7 & D6 & D5 & D4 & D3 & D2 & D1 & D0 \\
\midrule
80-8F & lecture & \multicolumn{8}{c}{données de la ROM de caractères} \\
80-8F & écriture & \multicolumn{8}{c}{adresse dans la ROM de caractères} \\
90-97 & lecture & FS & INT & -- & -- & B4 & B3 & B2 & B1 \\
90-97 & écriture & FSR & A/G & +1K & GM2 & GM1 & GM0 & CSS1 & CSS0 \\
98-9F & écriture & \multicolumn{8}{c}{adresse basse d'une cartouche d'E/S} \\
\bottomrule
\end{tabular}
\end{center}
En lecture : FS à l'état bas permet le défilement ; INT signale
l'interruption du curseur ; B1--B4 ne servent pas. En écriture : FSR autorise
l'interruption du curseur (0 : pas de curseur) ; A/G choisit le mode
graphique ; +1K change la page vidéo affichée ; GM0--GM2 choisissent le mode
graphique ; CSS0--CSS1 sélectionnent le jeu de couleurs.

\paragraph{Occupation de la RAM vidéo (fig.~10).} En modes 1 et 2, seuls
512 octets d'écran (\hex{8000}--\hex{81FF}) et 512 octets d'attributs
(\hex{9000}--\hex{91FF}) servent. En modes 3 et 4, l'écran occupe
\hex{8000}--\hex{8BFF}. Les modes 5 et 6 utilisent les 6\,Ko
(\hex{8000}--\hex{97FF}). Le signal +1K commute la page affichée sans changer
les adresses vues par le processeur, d'où des calculs d'adresse compliqués
dans l'interpréteur et la lenteur de la X-720. Les auteurs conseillent de se
placer dans un mode et d'écrire ses propres routines de tracé (voir LE PIÈGE
en version vidéo).

\paragraph{Codage des points.}
\begin{center}\small
\begin{tabular}{@{}lll@{}}
\toprule
Mode & Contenu & Codage \\
\midrule
1 & texte seul, 16 lignes $\times$ 32 colonnes & 1 octet par caractère (code ASCII), 512 octets \\
2 & texte ou graphique & 6 points par case et 2 bits de couleur \\
3 et 5 & graphique & 4 points par octet, 4 couleurs \\
4 et 6 & graphique & 8 points par octet, 1 bit par point \\
\bottomrule
\end{tabular}
\end{center}
Le mode 4 utilise 3\,Ko ; le mode 6 utilise 6\,Ko, avec une définition
doublée.

\subsection{Logiciel}
La ROM de la X-720 commence par le mot-clé \texttt{love}, détecté à la mise
sous tension (voir les \emph{Mystères}, §\,7.4). Le sous-programme
d'initialisation (\hex{A0CA}) redirige les crochets de la zone système vers
SCREEN, PAINT et COLOR, et étend PSET, PRESET, POINT, LINE et CIRCLE.

\begin{center}\small
\begin{tabular}{@{}l>{\ttfamily}l|l>{\ttfamily}l|l>{\ttfamily}l@{}}
\toprule
Instruction & \normalfont Entrée & Instruction & \normalfont Entrée & Instruction & \normalfont Entrée \\
\midrule
SCREEN & AAD2 & SCREEN(X,Y) & ABBD & COLOR & AA73 \\
PSET & ABFB & PRESET & ABF6 & POINT & AC14 \\
LINE & AC86 & PAINT & AD9B & CIRCLE & AED2 \\
\bottomrule
\end{tabular}
\end{center}

\begin{center}\small
\begin{tabular}{@{}>{\ttfamily}l>{\raggedright}p{4.6cm}|>{\ttfamily}l>{\raggedright\arraybackslash}p{6.2cm}@{}}
\toprule
\normalfont Adresse & Variable système & \normalfont Routine & Rôle \\
\midrule
00B6 & page active (0 = page 1) & A015 & LSETCUR : curseur (comme C0A1) \\
00B7 & page affichée & A050 & INTC : clignotement du curseur \\
00D1 & mode d'écran & A07C & ECRIRAM : écriture VRAM en cours ? \\
0015 & 0 caractères normaux, 1 caractères LCD & A08B & CONPUT : sortie écran \\
04D7 & caractères KANA & A0CA & initialisation de la X-720 \\
04E5 & couleur d'avant-plan & A2C2 & MODECR : changement de mode (A = 0 mode 1, 1 mode 2) \\
04E6 & couleur de fond & A355 & RETLCD : retour à l'écran LCD \\
04E7 & palette (0 ou 2) & A374 & VIDEO : affichage sur le téléviseur \\
04C6 & $X_i$, $Y_i$ de LINE & A4D1 & PUTCAR : envoi d'un caractère \\
04C8 & $X_f$, $Y_f$ de LINE & A9FB--AA73 & table des sous-programmes par mode \\
04CA & déplacement STEP & A631--A9FA & sous-programmes par mode \\
\bottomrule
\end{tabular}
\end{center}

\subsection{Quatre utilitaires \pg{80-82}}

\paragraph{Test de présence.} Écrit \hex{80} en \hex{8000} et relit ; sans
X-720, saut en SN ERROR (\hex{F1AA}).
\begin{verbatim}
10 REM TEST DE LA X-720
20 FOR I=&H1C00 TO &H1C0D
30 READ A$:POKE I,VAL("&H"+A$)
40 NEXT I
50 DATA 3E,80,32,00,80,3A,00,80,FE,80
60 DATA C2,AA,F1,C9
\end{verbatim}
(\texttt{LD A,\$80 / LD (\$8000),A / LD A,(\$8000) / CP \$80 / JP NZ,\$F1AA / RET})

\paragraph{Effacement partiel.} Fige les 8 premières lignes de l'écran
vidéo en réduisant la zone de défilement (\hex{00BB}--\hex{00BD}) et efface
le reste (\hex{CE9E}).
\begin{verbatim}
10 REM EFFACEMENT PARTIEL
20 FOR I=&H1C00 TO &H1C1D
30 READ A$:POKE I,VAL("&H"+A$):NEXT I
50 DATA CD,0D,1C,CD,9E,CE,21,08,01,22,B8,00,C9
60 DATA 3E,08,32,BB,00,3E,10,32,BC,00,3E,09,00
70 DATA 32,BD,00,C9
\end{verbatim}
\nota{Le livre écrit \texttt{LD A,(\&16-8+1)} et ses DATA donnent
\texttt{3A,10,00} (\texttt{LD A,(\$0010)}, lecture en mémoire) : il faut
\texttt{LD A,9} (\texttt{3E,09}). La ligne 60 est corrigée ici en
\texttt{3E,09,00}, où le \texttt{00} est un \texttt{NOP} qui garde les
adresses inchangées.}

\paragraph{\texttt{SCREEN}.} Le mode moins 1 va dans C et en \hex{00D1} ;
la page active moins 1 dans D, la page affichée moins 1 dans E ; puis saut en
\hex{AB09}.
\begin{verbatim}
10 REM *** SCREEN ***
20 FOR I=&H1C00 TO &H1C0C
30 READ A$:POKE I,VAL("&H"+A$):NEXT I
50 DATA 3E,00,32,D1,00,4F,3E,01,57
60 DATA 5F,C3,09,AB
\end{verbatim}
(\texttt{LD A,0 / LD (\$D1),A / LD C,A / LD A,1 / LD D,A / LD E,A /
JP \$AB09} : SCREEN 1, pages 2 et 2.)

\paragraph{\texttt{COLOR}.} Écrit directement les couleurs d'avant-plan, de
fond et la palette.
\begin{verbatim}
10 REM *** COLOR ***
20 FOR I=&H1C00 TO &H1C0F
30 READ A$:POKE I,VAL("&H"+A$):NEXT I
50 DATA 3E,01,32,E5,04,3E,02,32,E6,04
60 DATA 3E,00,32,E7,04,C9
\end{verbatim}

%======================================================================
\part*{3\textsuperscript{e} partie : Logiciel \pg{84-145}}
\addcontentsline{toc}{part}{3e partie : Logiciel}

%----------------------------------------------------------------------
\section{Présentation et chargement \pg{85}}

Chaque logiciel est présenté avec ses spécifications, son principe, son mode
d'emploi et son listing. Les logiciels en BASIC se chargent par
\texttt{CLOAD} depuis la cassette des Éditions Neptune. Les logiciels en
assembleur se chargent par \texttt{CLOAD} puis \texttt{RUN} (un chargeur
BASIC lit les codes) ; \texttt{RUN 100} fait une copie des codes sur
cassette.

Sans la cassette, chaque logiciel en assembleur est accompagné de deux
programmes types, à adapter avec ses adresses de début \texttt{D} et de fin
\texttt{F} :

\begin{verbatim}
5 REM *** ENTREUR DE CODES ***            10 CLEAR50,&H7FF
10 CLEAR50,&H7FF:A=&H800                  20 INIT#1,"CASI:"
20 PRINT HEX$(A);" : ";:INPUT C$          30 INPUT#1,N$,D,F:MOTOR
30 V=VAL("&H"+C$):POKE A,V                50 PRINT"Trouve :";N$
49 A=A+1:IF A>&HA29 THEN PRINT"...":      60 FORI=D-1TOF
   BEEP 2,3:END                           70 POKEI,INP(#1):NEXT:MOTOR:END
50 GOTO20                                 100 CLEAR50,&H7FF
                                          110 D=&H800:F=&HA29
                                          120 N$="LLIST":INIT#1,"CASO:"
                                          130 INPUT"Magneto OK";T$
                                          140 PRINT#1,N$,D,F:MOTOR
                                          150 FORI=1TO1800:NEXT
                                          160 FORI=DTOF:OUT#1,PEEK(I)
                                          170 NEXT:MOTOR:END
\end{verbatim}
L'entreur de codes (à gauche) fait taper les octets un par un, en
hexadécimal, à partir des vidages du livre. Le chargeur (à droite) relit les
codes sur cassette (lignes 10--90) ou les y enregistre (lignes 100--180).

La lecture commence en \texttt{D-1}, un octet plus tôt que l'écriture.
Ce n'est probablement pas une erreur : le premier octet lu est sans doute le
séparateur laissé par \texttt{PRINT\#1} après \texttt{F}, qui tombe ainsi
juste avant le programme. Les données arrivent donc bien de \texttt{D} à
\texttt{F}.

\subsection{Récapitulatif des logiciels}
\begin{center}\small
\begin{tabular}{@{}lll>{\ttfamily}l>{\ttfamily}l@{}}
\toprule
Logiciel & Type & Longueur & \normalfont Implantation & \normalfont Lancement \\
\midrule
LMDATA & BASIC & 1670 octets & -- & RUN \\
LOGO CANON & BASIC & 831 octets & -- & RUN \\
RESTORE calculé & BASIC + LM & 372 octets & 2000 & RUN \\
COPIE RAPIDE & BASIC + LM & 539 octets & chaîne & RUN \\
LOGOGÉNÈSE & BASIC + LM & 1819 octets & chaîne & RUN \\
LLIST & LM & 554 octets & 0800--0A29 (relogé) & EXEC \&H800 \\
LABYRINTHE 3D & LM & 2066 octets & 0800--1012 & EXEC \&H800 \\
SOLITAIRE & LM & 4609 octets & 0800--11D5 & EXEC \&H800 \\
PENTOMINOS & LM & 3489 octets & 0800--159A & EXEC \&H800 \\
AUTONUM & BASIC + LM & 1093 + 216 octets & relogeable & RUN \\
REFBAS & LM & 985 octets & 3C10--3FE9 & EXEC \&H3C10 \\
EXABAS & LM & 901 octets & 3C10--3F95 & EXEC \&H3C10 \\
LE PIÈGE (LCD) & LM & 4205 octets & 0800--186C & EXEC \&H800 \\
LE PIÈGE (vidéo) & LM & 4329 octets & 0800--18E8 & EXEC \&H800 \\
OTHELLO-REVERSI & LM & 3073 octets & 0800--1400 & EXEC \&HE00 \\
\bottomrule
\end{tabular}
\end{center}

%----------------------------------------------------------------------
\section{LMDATA et applications \pg{86-88}}

\paragraph{But.} Transcrire automatiquement une routine en mémoire sous
forme de lignes \texttt{DATA} décimales (trois chiffres par octet), à la
suite du programme. On évite ainsi de recopier les codes à la main.

\paragraph{Utilisation.} Implanter la routine, lancer LMDATA, donner les
adresses de début et de fin. Effacer ensuite les lignes DATA superflues, puis
les lignes 1 à 230 (commande DEL de la carte XP-140, ou à la main). Au-delà
de 256 octets, il faut ajouter des lignes DATA. Faire une copie de sécurité
avant.

\paragraph{Principe.} Les lignes 10000 à 10150 contiennent des DATA
« fictifs » (« \dots ») ; LMDATA y écrit directement, par \texttt{POKE}, les
codes ASCII des chiffres (\texttt{CA=48} est le code de « 0 »). La ligne 40
calcule l'adresse du premier champ de la ligne 10000 ; la ligne 210 décompose
chaque octet en centaines, dizaines et unités.
\begin{verbatim}
1 REM Ecriture automatique de codes LM en DATA
10 CLEAR50,&H1F00:DEFINTA-Z:LL=16:CA=48:CLS
20 INPUT"Adres. debut routine";DM:DD=DM
30 INPUT"Adresse fin routine";FM:NC=FM-DM+1:JC=INT(NC/LL)
40 RESTORE10000:AD=PEEK(&H328)+256*PEEK(&H329)+6
50 LB=LL:FORJ=1TOJC
55 PRINT"LIGNE";10000+10*J
60 GOSUB200:AD=AD+5:NEXTJ
70 LB=NC-JC*LL:IFLB=0THEN100
80 GOSUB200
90 PRINT"Effacer la fin de la ligne";10000+10*JC
100 PRINT"Effacer les lignes";10000+10*(JC+1);"";10150;
110 IFINKEY$=""THEN110
120 END
200 FORI=1TOLB:VC=PEEK(DM):DM=DM+1
210 C=INT(VC*.01):D=INT((VC-100*C)/10):U=VC-10*D-100*C
220 POKEAD,C+CA:POKEAD+1,D+CA:POKEAD+2,U+CA:AD=AD+4
230 NEXT:RETURN
10000 DATA...,...,...,...,...,...,...,...,...,...,...,...,...,...,...,...
  (idem jusqu'a la ligne 10150)
\end{verbatim}
\nota{La ligne 55 annonce la ligne $10000+10J$ alors que la ligne en cours
d'écriture est $10000+10(J-1)$ : l'affichage est décalé de dix.}

\nota{Vérifié dans l'émulateur (vraie ROM) : les lignes gabarits doivent être
tapées \textbf{sans espace} après \texttt{DATA} : \texttt{10000
DATA...,...,...}. Après \texttt{RESTORE}, le pointeur \hex{0328} désigne
l'octet qui précède la ligne : le \texttt{+6} de la ligne 40 tombe donc
juste après le jeton \texttt{DATA}. Le X-07 conserve l'espace tapée après
\texttt{DATA} ; avec \texttt{10000 DATA ...}, les chiffres sont écrits un
caractère trop tôt (\texttt{000.007.014}\dots, puis \texttt{12} au lieu de
\texttt{112} à la ligne suivante). C'est d'ailleurs sans espace que
LMDATA écrit les lignes données plus loin (\texttt{10000 DATA035,035,}\dots).}

\subsection*{Exemple 1 : \texttt{RESTORE} calculé (fig.~17)}
La routine du chapitre 1, implantée en \hex{2000} avec des DATA produits par
LMDATA.
\begin{verbatim}
1 REM Exemple de RESTORE CALCULE.
2 REM Cree par LMDATA
10 CLEAR50,&H1FFF:DEFINT A-Z
20 RESTORE10000
30 FORI=0TO17:READA:POKE&H2000+I,A:NEXT
40 INPUT"N= (1..3)";N
50 IFN<1 OR N>3 THEN 40
60 J=N*100
70 Z=USR(&H2000,J)
80 READA$:PRINTA$:GOTO40
100 DATA"lecture DATA 100"
200 DATA"lecture DATA 200"
300 DATA"lecture DATA 300"
999 END
10000 DATA035,035,094,035,086,205,013,243,096,105,210,056,246,043,034,040
10010 DATA003,201
\end{verbatim}

\subsection*{Exemple 2 : logo CANON (fig.~18)}
Programme BASIC qui dessine le logo CANON sur l'écran LCD par
\texttt{PSET}. Chaque ligne d'écran (6 à 25) est décrite par des longueurs de
segments, alternativement allumés et éteints, dont la somme vaut 79.
\begin{verbatim}
5 REM EXEMPLE DE GRAPHISME
10 CLEAR200:DEFINTA-Z
20 X=23:CLS
30 FORL=6TO25
40 A1=1:F=1
50 READA:A1=A1+A:IFF=1THEN70
60 F=1:GOTO80
70 FORI=A1-ATOA1-1:PSET(X+I,L):NEXT:F=0
80 IFA1=<79THEN50
90 NEXT
100 END
110 DATA79,9,4,66,8,6,65,6,10,63
120 DATA5,12,7,4,9,3,3,2,7,7,6,3,3,2,6
130 DATA4,6,6,2,5,6,7,5,1,4,5,9,4,5,1,4,5
140 DATA4,5,6,4,2,9,4,13,3,2,3,6,1,13,4
150 DATA3,6,6,2,10,3,6,5,2,5,2,3,2,6,3,5,2,5,3
160 DATA3,5,6,1,12,4,5,5,2,5,1,4,3,6,2,5,2,5,3
170 DATA3,5,15,8,5,5,2,5,1,4,3,6,2,5,2,5,3
180 DATA3,5,14,10,4,5,2,5,1,5,3,5,2,5,2,5,3
190 DATA3,5,13,5,2,4,4,5,2,5,1,5,3,5,2,5,2,5,3
200 DATA3,6,11,5,4,4,3,5,2,5,1,6,3,4,2,5,2,5,3
210 DATA4,5,9,1,1,5,4,4,3,5,2,5,1,6,3,4,2,5,2,5,3
220 DATA4,6,6,2,2,6,2,6,2,5,2,5,2,6,2,3,3,5,2,5,3
230 DATA5,12,4,8,1,4,2,5,2,5,2,6,3,2,3,5,2,5,3
240 DATA6,10,6,6,2,5,1,5,2,5,3,9,4,5,2,5,3
250 DATA8,6,9,4,3,5,1,5,2,5,4,7,5,5,2,5,3,79,79
\end{verbatim}
Les valeurs ont été relues sur le scan et contrôlées : chacune des 20 lignes
d'écran fait bien 79 points.

\subsection*{Exemple 3 : copie rapide de l'écran sur la X-710 (fig.~19)}
La routine (53 octets, relogeable) est implantée dans une chaîne. Elle
parcourt l'image miroir du LCD (\hex{0214}, 4 lignes de 20 caractères),
supprime les espaces de fin de ligne et envoie chaque ligne à l'imprimante
(\hex{CEF7}), suivie de CR et LF (\hex{CFB0}).
\begin{asm}
00 & 06 00    & & LD B,0         & numéro de ligne \\
02 & 21 14 02 & & LD HL,\$0214   & image miroir de l'écran \\
05 & C5       & L1 & PUSH BC     & \\
06 & E5       & & PUSH HL        & \\
07 & 01 13 00 & & LD BC,\$0013   & dernier caractère de la ligne \\
0A & 09       & & ADD HL,BC      & \\
0B & 7E       & L2 & LD A,(HL)   & \\
0C & FE 20    & & CP " "         & espace de fin ? \\
0E & 20 04    & & JR NZ,L3       & \\
10 & 2B       & & DEC HL         & \\
11 & 0D       & & DEC C          & \\
12 & 20 F7    & & JR NZ,L2       & \\
14 & 0C       & L3 & INC C       & C = nombre de caractères utiles \\
15 & E1       & & POP HL         & \\
16 & E5       & & PUSH HL        & \\
17 & 7E       & L4 & LD A,(HL)   & \\
18 & 04       & & INC B          & \\
19 & 23       & & INC HL         & \\
1A & E5       & & PUSH HL        & \\
1B & C5       & & PUSH BC        & \\
1C & CD F7 CE & & CALL \$CEF7    & impression du caractère \\
1F & C1       & & POP BC         & \\
20 & E1       & & POP HL         & \\
21 & 78       & & LD A,B         & \\
22 & B9       & & CP C           & \\
23 & 20 F2    & & JR NZ,L4       & \\
25 & CD B0 CF & & CALL \$CFB0    & CR + LF \\
28 & E1       & & POP HL         & \\
29 & 01 14 00 & & LD BC,\$0014   & ligne suivante \\
2C & 09       & & ADD HL,BC      & \\
2D & C1       & & POP BC         & \\
2E & 04       & & INC B          & \\
2F & 78       & & LD A,B         & \\
30 & FE 04    & & CP 4           & 4 lignes ? \\
32 & 20 D1    & & JR NZ,L1       & \\
34 & C9       & & RET            & \\
\end{asm}
\begin{verbatim}
1 REM Copie ecran
2 REM Methode de la chaine
3 REM DATA crees par LMDATA.
10 CLEAR 100:DEFINT A-Z
20 CH$=STRING$(53,"#")
30 GOSUB 300
40 FORI=0TO52:READA:POKEAD+I,A:NEXT
100 PRINT" Copie rapide de l'"
110 PRINT"      ecran sur"
120 PRINT" l'imprimante X-710"
130 PRINT"      +++";
200 GOSUB300:EXEC AD
210 END
300 AD=VARPTR(CH$):AD=PEEK(AD+1)+256*PEEK(AD+2):RETURN
10000 DATA006,000,033,020,002,197,229,001,019,000,009,126,254,032,032,004
10010 DATA043,013,032,247,012,225,229,126,004,035,229,197,205,247,206,193
10020 DATA225,120,185,032,242,205,176,207,225,001,020,000,009,193,004,120
10030 DATA254,004,032,209,201
\end{verbatim}

%----------------------------------------------------------------------
\section{LOGOGÉNÈSE \pg{89-90}}\label{sec:logo}

\paragraph{But.} Créer des mots savants au hasard en associant un préfixe
grec ou latin (« HYDRO », « PHILO »\dots) et un suffixe (« LOGIE »,
« PHOBE »\dots). Exemples possibles : « CHRONOPHAGE », « XYLOPHOBE ».

\paragraph{Principe.} Le programme BASIC (1819 octets) tire une ligne de
préfixes (100 à 132) et une ligne de suffixes (200 à 213). Il s'y place grâce
au \texttt{RESTORE} calculé du chapitre 1, puis lit un mot au hasard (1 à 4)
dans chaque ligne. La routine est logée dans la chaîne \texttt{RM\$} ; son
adresse est recalculée avant chaque appel (ligne 90), car la zone des
chaînes se déplace.
\begin{verbatim}
1 'ATTENTION NE PAS RENUMEROTER
10 DEFINTA-Z:CLS:PRINT"Logogenese"
20 RM$=STRING$(18,0):GOSUB90
30 FORI=0TO17:READB$:POKEAD+I,VAL("&H"+B$):NEXT:Z=RND(0)
40 DATA23,23,5E,23,56,CD,0D,F3,60,69,D2,38,F6,2B,22,28,03,C9
45 J=INT(RND(1)*33)+100:GOSUB90
50 Z=USR(AD,J):N=INT(RND(1)*4+1)
55 FORI=1TON:READD$:NEXT
60 K=INT(RND(1)*14)+200:GOSUB90
65 Z=USR(AD,K):N=INT(RND(1)*4+1)
70 FORI=1TON:READF$:NEXT
75 PRINTD$;F$:IFINKEY$=""THEN45
80 IFINKEY$=""THEN80ELSE45
90 AD=VARPTR(RM$):AD=PEEK(AD+1)+256*PEEK(AD+2):RETURN
95 END
\end{verbatim}

{\small
\begin{longtable}{@{}>{\ttfamily}l>{\ttfamily}l|>{\ttfamily}l>{\ttfamily}l@{}}
\toprule
\normalfont Ligne & \normalfont Préfixes & \normalfont Ligne & \normalfont Préfixes / suffixes \\
\midrule\endhead
\bottomrule\endfoot
100 & AERO,AGRO,AMBI,ANDRO & 117 & MORPHO,MYTHO,NECRO,NEO \\
101 & ANTI,ASTRO,BIBLIO,BIO & 118 & NOSO,NUCLEO,OCTO,OLEO \\
102 & CACO,CALLI,CLEPTO,CHIMIO & 119 & OMNI,ORTHO,PALEO,PAN \\
103 & CHRONO,CINE,COPRO,COSMO & 120 & PARA,PAPYRO,PATHO,PEDO \\
104 & CRYO,CRYPTO,DECI,DECA & 121 & PENTA,PERI,PETRO,PHAGO \\
105 & HEMI,DERMO,DEXTRO,DODECA & 122 & PHALLO,PHANERO,PHILO,PHYTO \\
106 & DROMO,DYNAMO,DYS,ECTO & 123 & PHOBO,PHONO,PHOTO,PHRENO \\
107 & ELECTRO,EMBRYO,ENTOMO & 124 & PHYSIO,POLY,MONO,HEMI \\
108 & ERGO,EROTICO,GASTRO,GEO & 125 & PORNO,POST,PROTO,PSEUDO \\
109 & GLOSSO,GONIO,GONO,GYMNO & 126 & PSYCHO,RADIO,RETRO,RHINO \\
110 & HECTO,HELIO,HEMATO,HETERO & 127 & SADO,STEREO,SCLERO,SEMI \\
111 & HIERO,HIPPO,HOLO,HOMEO & 128 & SCHIZO,SIMILI,SPELEO,STENO \\
112 & HOMO,HYDRO,HYPER,HYPO & 129 & SUB,SUPRA,SUPER,TACHY \\
113 & ICONO,IDEO,INFRA,ISO & 130 & TECHNO,TELE,TELEO,THEO \\
114 & INTRA,LATERO,LIPO,LOGO & 131 & THERMO,TRIBO,ULTRA,VIDEO \\
115 & LOXO,MACRO,MEGALO,METEO & 132 & XENO,XYLO,ZOO,MINI \\
116 & META,MICRO,MISO,MNEMO & & \\
\midrule
200 & ALGIE,CARDE,CINESE,CEPHALE & 207 & MATIQUE,MEGALIE,METRIE,MNESIE \\
201 & CLASTE,COSMOS,CYCLE,DACTYLE & 208 & MORPHE,NAUTE,NEVROSE,NOMIE \\
202 & DERME,DIDACTE,DOXE,DROME & 209 & PATHE,PHAGE,PHANIE,PHILE \\
203 & DYNE,FUGE,GAME,GASTRE & 210 & PHOBE,PHONE,PNEE,PYGE \\
204 & GLOTTE,GENESE,GENE,GRADE & 211 & RRHEE,SCAPHE,SCOPE,SCOPIE \\
205 & GYNE,LATERE,LATRE,LOGIE & 212 & STASE,STENIE,THEISME,TONIQUE \\
206 & LOGUE,MANCIE,MANE,MANIE & 213 & TROPE,TYPE,VORE,TIQUE \\
\end{longtable}
}
Chaque ligne commence par \texttt{DATA} (par exemple
\texttt{100 DATA AERO,AGRO,AMBI,ANDRO}).

\nota{La ligne 107 n'a que trois mots : un tirage de 4 lit le premier mot de
la ligne 108. C'est sans conséquence. La ligne 116 est peu lisible sur le
scan (« JICRO, MKsO ») ; « MICRO, MISO » est la lecture probable.}

%----------------------------------------------------------------------
\section{LLIST \pg{91-95}}

\paragraph{But.} Améliorer la présentation des listings BASIC à
l'imprimante : titre, numérotation des pages, saut de page, alignement des
lignes et formatage des numéros. La figure~23 montre le résultat sur le logo
CANON (« LISTING LOGO CANON -- Page: 1 »).

\paragraph{Utilisation.}
\begin{enumerate}
  \item Vérifier par \texttt{DIR} et \texttt{FSET} qu'il reste au moins
        600 octets libres dans la zone des fichiers RAM : la partie active
        s'installe au fond de la mémoire.
  \item Charger les codes (\hex{0800}--\hex{0A29}).
  \item Régler la longueur des pages (octet \hex{0950}) et des lignes (octet
        \hex{093C}, qui contient \hex{27}). Mettre \hex{26} évite des sauts
        de ligne imprévus avec une ligne de 39 caractères.
  \item \texttt{EXEC \&H800} : le relogeur copie LLIST au fond de la
        mémoire, affiche « NOTEZ AD= \dots » puis se détruit. Il faut noter
        cette adresse (décimale, variable selon la version du X-07) : LLIST
        se lance ensuite par \texttt{EXEC adresse}.
\end{enumerate}
LLIST analyse le programme comme REFBAS et EXABAS (voir
§\,\ref{sec:refbas}). Il utilise la routine ROM \hex{FEFE}, qui recopie dans
le tampon d'entrée (\hex{00D5}) la chaîne pointée par HL (terminée par 0).
Longueur : 554 octets.

\subsection*{Le relogeur}
Ce petit programme peut servir à reloger n'importe quelle routine. Il
utilise une table des adresses à corriger, construite par l'assembleur (qui
doit savoir calculer sur les étiquettes) ou à la main. Il ne faut en oublier
aucune, et la table est à refaire à chaque modification du logiciel.
\begin{src}
DEB & LD HL,(\$212) & RAMEND : fond de la mémoire \\
 & DEC HL & \\
 & LD BC,fin-début & longueur à déplacer \\
 & XOR A & CY = 0 \\
 & SBC HL,BC & HL = nouvelle adresse de début \\
 & PUSH HL & (3 fois) \\
 & CALL \$CE9E & efface l'écran \\
 & LD HL,MSG1 & « NOTEZ AD = » \\
 & CALL \$FEF7 & \\
 & POP HL & \\
 & CALL \$BB98 & affichage décimal de l'adresse \\
 & LD HL,MSG2 & « Pressez une touche » \\
 & CALL \$FEF7 & \\
KEY & XOR A & \\
 & CALL \$C90A & attente d'une touche \\
 & JR Z,KEY & \\
 & POP HL & début après transfert \\
 & LD DE,début & \\
 & XOR A & \\
 & SBC HL,DE & HL = décalage \\
 & PUSH HL & \\
 & POP BC & BC = décalage \\
 & LD IX,table & adresses à modifier \\
DEB0 & LD L,(IX+0) & \\
 & LD H,(IX+1) & HL = adresse à modifier \\
 & LD E,(HL) & \\
 & INC HL & \\
 & LD D,(HL) & DE = contenu \\
 & DEC HL & \\
 & PUSH HL & \\
 & EX DE,HL & \\
 & ADD HL,BC & contenu + décalage \\
 & EX DE,HL & \\
 & POP HL & \\
 & LD (HL),E & modification \\
 & INC HL & \\
 & LD (HL),D & \\
 & INC IX & entrée suivante \\
 & INC IX & \\
 & LD A,(IX+0) & \\
 & CP 0 & fin de la table ? \\
 & JR NZ,DEB0 & \\
 & POP DE & DE = début (mémoire haute) \\
 & LD HL,début & HL = début (mémoire basse) \\
 & LD BC,fin-début & \\
 & LDIR & transfert du programme modifié \\
 & LD HL,DEB & \\
 & LD A,\$C3 & écrit C3 C3 C3 sur DEB : \\
 & LD B,3 & le relogeur devient inutilisable \\
DEB1 & LD (HL),A & \\
 & INC HL & \\
 & DJNZ DEB1 & \\
 & JP \$C3C3 & retour au BASIC \\
MSG1 & DB 'NOTEZ AD =',0 & \\
MSG2 & DB 'Pressez ...',0 & \\
TABLE & DW A1+1, \dots & une entrée par adresse absolue \\
 & DB 0 & fin de la table \\
\end{src}
\nota{La fin de la table est testée sur l'octet de poids faible seulement :
une adresse à corriger dont l'octet faible vaut \hex{00} arrêterait la
relocation. Le vidage de LLIST (fig.~21, \hex{0800}--\hex{0A29}) n'est pas
recopié ici.}

%----------------------------------------------------------------------
\section{LABYRINTHE 3D \pg{96-100}}

\paragraph{But.} Sortir le plus vite possible d'un labyrinthe. Deux vues
sont possibles : de « haut » (le dédale entier) ou de « bas » (vue en
perspective de l'endroit où l'on se trouve). Une boussole guide le joueur et
un score (le temps) s'affiche : le plus bas est le meilleur.

\paragraph{Utilisation.} \texttt{EXEC \&H800}, puis une touche après la
présentation. Commandes :
\begin{center}\small
\begin{tabular}{@{}cl|cl@{}}
\toprule
N, W, E, S & tourner vers le nord, l'ouest, l'est, le sud & A & avancer \\
L & vue d'ensemble (coûte 20 points) & R & repartir du départ \\
B & retour au BASIC & & \\
\bottomrule
\end{tabular}
\end{center}
Le labyrinthe est créé automatiquement ; l'entrée est d'un côté, la sortie
à l'opposé.

\paragraph{Algorithme de création du dédale} (exploration en profondeur avec
retour arrière) :
\begin{enumerate}
  \item créer un tableau de cellules vides ;
  \item choisir une origine proche du centre ;
  \item tester les quatre cellules voisines pour savoir lesquelles ont été
        visitées ;
  \item si toutes l'ont été, revenir en arrière jusqu'à une cellule visitée
        voisine d'une cellule non visitée ;
  \item choisir au hasard une voisine non visitée ;
  \item y aller en ouvrant une porte dans les deux cellules, et marquer la
        nouvelle cellule ;
  \item tant que toutes les cellules n'ont pas été testées, revenir en 3 ;
        sinon créer une entrée et une sortie.
\end{enumerate}
Longueur : 2066 octets, de \hex{0800} à \hex{1012} (vidage fig.~24). Pour
les chargeurs (fig.~25) : \texttt{CLEAR50,\&H7F0}, \texttt{D=\&H800},
\texttt{F=\&H1012}, nom de fichier « PLABIR ».

\nota{Le texte annonce 20 points pour la vue d'ensemble ; le message stocké
dans le programme semble indiquer 22 points.}

%----------------------------------------------------------------------
\section{Le SOLITAIRE \pg{101-106}}

\paragraph{Le jeu.} Le solitaire se joue sur une planchette de 33, 37 ou 41
trous ; le logiciel utilise la version anglaise à 33 trous. Toutes les
fiches sont en place sauf une. Une fiche saute horizontalement ou
verticalement par-dessus une fiche voisine pour tomber dans un trou libre,
et la fiche sautée est retirée. La partie s'arrête quand plus aucun coup
n'est possible ; l'idéal est de ne garder qu'une fiche. Il n'y a aucun
hasard.

\begin{center}\small\ttfamily
\begin{tabular}{@{}*{7}{c}@{}}
   &    &  1 &  2 &  3 &    &    \\
   &    &  4 &  5 &  6 &    &    \\
 7 &  8 &  9 & 10 & 11 & 12 & 13 \\
14 & 15 & 16 & 17 & 18 & 19 & 20 \\
21 & 22 & 23 & 24 & 25 & 26 & 27 \\
   &    & 28 & 29 & 30 &    &    \\
   &    & 31 & 32 & 33 &    &    \\
\end{tabular}
\end{center}

\paragraph{Utilisation.} \texttt{EXEC \&H800}.
\begin{itemize}
  \item \textbf{Option 1, recherche par le X-07.} On choisit le nombre de
        fiches de la solution finale et l'un des 22 ordres d'examen des
        directions. Il n'existe pas d'algorithme simple : le programme
        essaie systématiquement les coups et revient en arrière niveau par
        niveau en cas d'impasse. L'ordre d'examen joue donc beaucoup sur la
        durée : avec 3 fiches ou plus, quelques secondes ; avec 2 fiches,
        de 5 minutes à une heure ; avec l'ordre 9, une solution classique
        en moins de 10 minutes.
  \item \textbf{Option 2, partie du joueur.} Le curseur se déplace avec les
        flèches. « D » marque la case de départ, « A » la case d'arrivée et
        joue le coup, « - » annule le coup précédent, « Q » arrête quand la
        partie est bloquée, CTRL Q revient au BASIC.
\end{itemize}
L'affichage de la profondeur d'analyse ralentit beaucoup la recherche ; on
le supprime en mettant à 0 les octets \hex{0989}, \hex{098A} et \hex{098B}
par \texttt{POKE}.

Longueur : 4609 octets, de \hex{0800} à \hex{11D5} (vidage fig.~26).
Chargeurs (fig.~27) : \texttt{D=\&H800}, \texttt{F=\&H11D5}, nom « SOL ».

%----------------------------------------------------------------------
\section{Les PENTOMINOS \pg{107-113}}

\paragraph{Le puzzle.} Les 12 pentominos sont les 12 façons d'assembler
cinq carrés par leurs côtés. On les désigne par une lettre de forme
voisine : F, I, L, N, P, S, T, U, V, W, X, Y. Le casse-tête est connu depuis
1907 (\emph{Canterbury Puzzles}). Nombre de façons de remplir un rectangle
avec les 12 pièces :
\begin{center}\small
\begin{tabular}{@{}lcccc@{}}
\toprule
Rectangle & $3\times20$ & $4\times15$ & $5\times12$ & $6\times10$ \\
Solutions & 2 & 368 & 1010 & 2339 \\
\bottomrule
\end{tabular}
\end{center}

\paragraph{Utilisation.} \texttt{EXEC \&H800}, une touche après chaque
affichage, puis choix du rectangle et de l'ordre d'examen des pièces. Le
programme cherche par essais et erreurs ; l'ordre conditionne la durée et la
solution trouvée. Le résultat peut être imprimé sur la X-710. CTRL Q revient
au BASIC. L'ordre « TFILUVNWPXSY » donne une solution rapide pour les trois
rectangles.

Longueur : 3489 octets, de \hex{0800} à \hex{159A} (vidage fig.~28).
Chargeurs (fig.~29) : \texttt{D=\&H800}, \texttt{F=\&H159A}, nom « PENTA ».

%----------------------------------------------------------------------
\section{AUTONUM \pg{114-116}}

\paragraph{But.} Ajouter au X-07 la numérotation automatique des lignes
BASIC (fonction AUTO).

\paragraph{Installation.} Choisir l'adresse de fin de la routine selon la
mémoire du X-07 et la protéger par \texttt{FSET} (215 octets au moins).
\texttt{RUN} demande cette adresse (en hexadécimal), affiche l'adresse de
début (à noter) et lance la routine. Le chargeur peut ensuite être effacé
par \texttt{NEW}.

\paragraph{Utilisation.}
\begin{itemize}
  \item CTRL N : un « ? » s'affiche ; taper le premier numéro et le pas, par
        exemple \texttt{100,10} : les lignes seront numérotées 100, 110,
        120\dots
  \item CTRL Z : active ou suspend la numérotation.
  \item CTRL Q : désactive AUTONUM. Pour le relancer : \texttt{EXEC AD}, avec
        l'adresse affichée à l'installation.
\end{itemize}

\paragraph{Principe.} AUTONUM détourne l'interruption du clavier (INT A,
vecteur en \hex{003D}). À chaque RETURN, la ligne est d'abord traitée
normalement ; si AUTONUM est actif, le numéro suivant est ensuite écrit dans
le tampon du clavier. Les cinq octets de données sont rangés dans le tampon
de CSAVE/CLOAD : \hex{0300} (indicateur actif/inactif), \hex{0301}--\hex{0302}
(numéro courant), \hex{0303}--\hex{0304} (pas).

La routine est relogeable : elle doit trouver sa propre adresse pour
l'écrire dans le vecteur. L'astuce consiste à écrire en \hex{0300} la
séquence \texttt{POP DE / PUSH DE / RET}, puis à l'appeler. Au retour, DE
contient l'adresse de l'instruction qui suit le \texttt{CALL}.

\begin{asm}
00 & 21 D1 D5 & & LD HL,\$D5D1   & D1 = POP DE, D5 = PUSH DE \\
03 & 22 00 03 & & LD (\$0300),HL & \\
06 & 3E C9    & & LD A,\$C9      & RET \\
08 & 32 02 03 & & LD (\$0302),A  & \\
0B & CD 00 03 & & CALL \$0300    & empile l'adresse XXX \\
0E & 21 0E 00 & XXX & LD HL,14   & au retour, DE = XXX \\
11 & 19       & & ADD HL,DE      & HL = AUTO (14 octets plus loin) \\
12 & F3       & & DI             & \\
13 & 22 3D 00 & & LD (\$003D),HL & nouveau vecteur de INT A \\
16 & FB       & & EI             & \\
17 & AF       & & XOR A          & \\
18 & 32 00 03 & & LD (\$0300),A  & AUTONUM inactif au départ \\
1B & C9       & & RET            & \\
1C & D9 08    & AUTO & EXX ; EX AF,AF' & traitement de l'interruption \\
1E & DB F2 \dots & & IN A,(\$F2)\dots & tri des sources (comme la ROM) \\
33 & DB F1    & & IN A,(\$F1)    & touche reçue du T6834 \\
35 & FE 0D    & & CP \$0D        & RETURN : numéro suivant \\
84 & FE 1A    & & CP \$1A        & CTRL Z : bascule \\
A3 & FE 0E    & & CP \$0E        & CTRL N : saisie du numéro et du pas \\
D0 & FE 11    & & CP \$11        & CTRL Q \\
D2 & C2 80 C8 & & JP NZ,\$C880  & autre touche : traitement normal \\
D5 & C3 C3 C3 & & JP \$C3C3      & réinitialisation (vecteur d'origine) \\
\end{asm}
\nota{Le texte du livre écrit \texttt{LD (\$300),(HL)} (instruction
inexistante) : c'est \texttt{LD (\$300),HL}. Il parle aussi de « cinq octets
de mémoire morte » ; il s'agit de RAM. La routine fait 216 octets : le
chargeur écrit de \texttt{D=F-215} à \texttt{F} inclus.}

\begin{verbatim}
1 REM fonction AUTONUM pour X-07
10 CLEAR50,&HFFFF:DEFINTA-Z:CLS
20 INPUT"Adr. fin ";A$:F=VAL("&H"+A$)
30 D=F-215:PRINT"Adr. deb = ";HEX$(D)
40 FORI=DTOF:READA:POKEI,A:NEXT
50 EXECD:END
10000 DATA033,209,213,034,000,003,062,201,050,002,003,205,000,003,033,014
10010 DATA000,025,243,034,061,000,251,175,050,000,003,201,217,008,219,242
10020 DATA230,001,202,193,200,219,240,230,192,040,008,230,128,202,053,200
10030 DATA195,012,200,219,241,254,013,032,075,095,175,205,098,194,175,205
10040 DATA170,194,058,000,003,183,202,189,200,062,001,211,245,217,008,197
10050 DATA213,229,245,237,091,001,003,042,003,003,025,034,001,003,235,205
10060 DATA156,187,035,205,110,213,205,004,215,205,059,202,010,197,213,095
10070 DATA175,205,098,194,175,205,170,194,209,193,021,003,032,238,241,225
10080 DATA209,193,251,201,254,026,032,027,062,001,211,245,217,008,197,213
10090 DATA229,245,062,013,239,062,010,239,033,000,003,175,182,047,119,032
10100 DATA221,024,176,254,014,032,041,062,001,211,245,217,008,197,213,229
10110 DATA245,175,050,000,003,205,242,235,215,205,204,255,237,083,001,003
10120 DATA207,044,205,204,255,237,083,003,003,175,061,050,000,003,024,131
10130 DATA254,017,194,128,200,195,195,195
\end{verbatim}
Ces lignes sont regroupées par 16 valeurs au lieu de la mise en page du
livre ; les valeurs sont celles du livre, et le désassemblage du programme
complet est cohérent.

%----------------------------------------------------------------------
\section{REFBAS \pg{117-120}}\label{sec:refbas}

\paragraph{But.} Obtenir la liste triée des variables d'un programme BASIC
avec les lignes où elles apparaissent, ou la table des références croisées
des numéros de ligne. Utile pour analyser, renommer, réorganiser ou
documenter un long programme.

\paragraph{Utilisation.} Avant de charger REFBAS : \texttt{FSET 1024}
(\texttt{FSET 4096} si la carte moniteur est en place), pour le protéger en
fin de mémoire (version 16\,Ko). Lancement : \texttt{EXEC \&H3C10}.
\begin{itemize}
  \item Option V : table triée des variables et des lignes où elles sont
        utilisées.
  \item Option S : références croisées pour GOTO, GOSUB, THEN, ELSE,
        RESUME, RUN et RESTORE. Chaque ligne cible est suivie des lignes qui
        y renvoient.
  \item Option B : retour au BASIC (CTRL Q fonctionne à tout moment ; BREAK
        pendant l'impression).
\end{itemize}
REFBAS construit, à la suite du programme analysé, une table de 6 octets par
variable ou référence. Quand elle atteint la pile, il annonce « Mémoire
insuffisante » : il faut alors analyser le programme par morceaux. La longueur
des lignes imprimées (40 pour la X-710) se règle à l'adresse \hex{3EDB}.

\paragraph{Exemple (sur le chargeur de la fig.~32).}
\begin{verbatim}
Liste des variables
D      30 60 110 140 160
F      30 60 110 140 160
I      60 70 150 160 160
N$     30 50 120 140
T$     130
\end{verbatim}

\paragraph{Parcours d'un programme BASIC.} Comme tous les BASIC Microsoft,
une ligne se compose de 2 octets de chaînage (adresse de la ligne suivante),
2 octets de numéro, du code (mots-clés sur un octet, le reste en ASCII) et
d'un octet 0. Le début du programme est pointé par TXTTAB (\hex{00B2}),
normalement \hex{0553}. REFBAS, EXABAS et LLIST utilisent la même boucle :
\begin{src}
 & LD IX,(TXTTAB) & \\
BOUCLE & XOR A & \\
 & CP (IX+1) & chaînage nul : fin du programme \\
 & JP Z,fin & \\
 & PUSH IX & début de la ligne \\
 & LD L,(IX+0) & \\
 & LD H,(IX+1) & HL = ligne suivante \\
 & LD E,(IX+2) & \\
 & LD D,(IX+3) & DE = numéro de ligne \\
 & PUSH HL & \\
 & POP IX & IX = ligne suivante \\
 & POP HL & \\
 & LD BC,4 & \\
 & ADD HL,BC & HL = premier code de la ligne \\
 & \dots & traitement \\
 & JP BOUCLE & \\
\end{src}
Longueur : 985 octets, de \hex{3C10} à \hex{3FE9} (vidage fig.~31).
Chargeurs : \texttt{CLEAR50,\&H9FF}, \texttt{D=\&H3C10}, \texttt{F=\&H3FE9},
nom « REFBAS ».

%----------------------------------------------------------------------
\section{EXABAS \pg{121-124}}

\paragraph{But.} Compléter REFBAS par une recherche d'occurrences et un
« chercher-remplacer » (fonctions REF et SR des grands systèmes).

\paragraph{Utilisation.} \texttt{FSET 1024} (ou 4096 avec la carte
moniteur), chargement comme REFBAS, puis \texttt{EXEC \&H3C10}. Menu :
\begin{center}\small
\begin{tabular}{@{}cl|cl@{}}
\toprule
V & recherche d'une variable & c & remplacement d'une chaîne \\
C & recherche d'une chaîne (hors mots-clés) & t & remplacement d'un mot-clé \\
T & recherche d'un mot-clé & B & retour au BASIC \\
\bottomrule
\end{tabular}
\end{center}
EXABAS affiche les numéros des lignes trouvées ou modifiées. Il marque une
pause quand l'écran est plein ; la touche « I » envoie alors l'écran à la
X-710. Un carreau signale la fin de l'analyse ; CTRL Q revient au BASIC.
Faites une copie du programme avant un remplacement : la longueur des lignes
change et EXABAS doit recalculer tous les chaînages qui suivent.

\paragraph{Exemple (sur le chargeur de la fig.~34).}
\begin{verbatim}
Nom? N            30 50 120 140
Mot cl? PRINT     50 140
Mot cl? FOR       60 150 160
Mot cl? DEFSTR
Nom? T$           130
\end{verbatim}
Longueur : 901 octets, de \hex{3C10} à \hex{3F95} (vidage fig.~33).
Chargeurs : \texttt{CLEAR50,\&H9FF}, \texttt{D=\&H3C10}, \texttt{F=\&H3F95},
nom « EXABAS ».

%----------------------------------------------------------------------
\section{LE PIÈGE \pg{125-137}}

\paragraph{But.} Un jeu d'aventures, genre alors presque inexistant sur le
X-07 : explorer une grande maison. Il faut visiter toutes les salles pour
atteindre la sortie.

\paragraph{Utilisation.} Deux versions, LCD et vidéo (X-720). Faire un
\texttt{FSET} avant le chargement : la mémoire de \hex{0800} à \hex{1900}
est utilisée.
\begin{center}\small
\begin{tabular}{@{}>{\ttfamily}ll@{}}
\toprule
\normalfont Lancement & Effet \\
\midrule
EXEC \&H800 & début du jeu \\
EXEC \&H815 & reprise à l'endroit quitté par « P » \\
EXEC \&H838 & reprise à l'endroit quitté par « Q » \\
\bottomrule
\end{tabular}
\end{center}
L'écran affiche : ligne 1, le nom de la salle et les directions possibles ;
ligne 2, les objets présents ; ligne 3, une description ; ligne 4, les
commandes et les réponses.

\paragraph{Commandes.} Une lettre : N, S, E, W pour se déplacer, I pour
l'inventaire, Q pour quitter, P pour une pause (sauvegarde). Sinon, une
phrase « verbe à l'infinitif + nom », par exemple
\texttt{SAUTER COULOIR}. Seules les 4 premières lettres comptent
(\texttt{SAUT COUL} suffit). Un verbe inconnu affiche « ???? », un nom
inconnu «     ???? ».

\begin{center}\small
\begin{tabular}{@{}lccl@{}}
\toprule
Version & Longueur & Implantation & Chargeurs \\
\midrule
LCD (fig.~35--36) & 4205 octets & \hex{0800}--\hex{186C} & \\
Vidéo (fig.~37--38) & 4329 octets & \hex{0800}--\hex{18E8} & \texttt{CLEAR50,\&H7F0}, \texttt{F=\&H18E8}, nom « AVENT » \\
\bottomrule
\end{tabular}
\end{center}

%----------------------------------------------------------------------
\section{OTHELLO-REVERSI \pg{138-144}}

\paragraph{Le jeu.} Le Reversi a été inventé en Angleterre dans la seconde
moitié du XIX\textsuperscript{e} siècle, puis « réinventé » au Japon sous le
nom d'Othello. Il se joue sur un plateau de $8\times8$ avec des pions blancs
d'un côté et noirs de l'autre ; quatre pions sont posés au centre au départ.
Chacun à son tour pose un pion de sa couleur de façon à encadrer au moins un
pion adverse (en ligne, en colonne ou en diagonale) ; les pions encadrés sont
retournés. Un retournement n'en provoque pas d'autre. Un joueur qui ne peut
pas retourner de pion passe son tour. La partie s'arrête quand plus personne
ne peut jouer ; le gagnant a le plus de pions. Sur le X-07, les deux camps
sont représentés par des pions de formes différentes.

\paragraph{Utilisation.} \texttt{EXEC \&HE00} puis RETURN. Après la
présentation, on choisit de jouer contre le X-07 ou contre un humain, le
niveau (1 facile à 9 difficile) et ses initiales (4 lettres). Le damier
s'affiche à gauche, le nombre de pions de chaque joueur à droite. Le point
clignotant se déplace avec les quatre flèches ; RETURN joue le coup. « P »
permet de passer si l'on est vraiment bloqué ; sinon, le X-07 affiche les
coups possibles. CTRL Q quitte le jeu. Les auteurs jugent le programme très
fort.

Longueur : 3073 octets, de \hex{0800} à \hex{1400} ; point d'entrée
\hex{0E00}.

%----------------------------------------------------------------------
\section{Conclusion du livre \pg{145}}
André Tonic rappelle que les deux livres sur le X-07 ont été écrits en six
mois. Le service technique du Club C7 répond aux questions sur les deux
ouvrages (courrier, téléphone, Minitel). Les Éditions Neptune promettent
d'autres livres à bas prix. Le texte se termine sur « Bravo au X-07 ».

%======================================================================
\part*{Annexes \pg{146-155}}
\addcontentsline{toc}{part}{Annexes}

\section{Assembleurs et désassembleurs (annexe 1)}
\begin{center}\small
\begin{tabular}{@{}>{\raggedright}p{4.2cm}>{\raggedright\arraybackslash}p{12.3cm}@{}}
\toprule
Logiciel & Appréciation des auteurs \\
\midrule
Assembleur Logi'Stick & correct, mais étiquettes limitées, pas de directives (ORG, DEFM\dots), emploi compliqué \\
Assembleur Micro Systèmes & publié dans le n\textsuperscript{o}\,49 ; écrit en langage machine, tient sur une carte de 4\,Ko, nombreuses pseudo-instructions ; préféré des auteurs, qui l'ont utilisé pour tous les listings du livre \\
Désassembleur-moniteur Logi'Stick & sur la même cassette que l'assembleur ; très bien fait ; 125\,F au Club C7 (adhérents) \\
Carte Canon XP-140F & 8\,Ko de ROM et 4\,Ko de RAM ; 10 instructions BASIC en plus (RENUM, MERGE, AUTO\dots) et un moniteur complet (points d'arrêt, désassemblage, entrée de codes) ; 500 à 550\,F \\
Désassembleur-moniteur Micro Systèmes & publié dans le n\textsuperscript{o}\,42 ; complet, un peu moins fouillé que celui de Logi'Stick \\
\bottomrule
\end{tabular}
\end{center}

\section{Le Club C7 (annexe 2)}
Association loi 1901 fondée en octobre 1984 par des professionnels de
l'informatique et des étudiants. Elle propose une gazette bimestrielle,
\emph{Le Son du Canon} (environ 40 pages A4, envoyée avec une cassette des
programmes), une programmathèque, une aide technique, une coopérative et des
stages. Cotisation annuelle : 380\,F (430\,F pour l'étranger).

\section{Les mystères du X-07 (annexe 3)}
Présentation du premier livre : 160 pages en trois parties (assembleur Z80,
mystères du X-07 avec près de 150 adresses, huit applications) et plus de
40 schémas.

\section{La cassette (annexe 4)}
Une cassette des Éditions Neptune (90\,F, plus 15\,F de port) contient les
15 programmes du livre : LMDATA, LOGO CANON, RESTORE, COPIE RAPIDE, LLIST,
LOGOGÉNÈSE, LABYRINTHE, PENTOMINOS, SOLITAIRE, AUTONUM, REFBAS 16\,Ko,
EXABAS 16\,Ko, LE PIÈGE (LCD et vidéo) et OTHELLO. Les programmes en
assembleur y sont fournis en BASIC et en codes (B+C).

\section{Bibliographie (annexe 5)}
\begin{center}\small
\begin{tabular}{@{}ll@{}}
\toprule
Éditeur & Titre \\
\midrule
Sybex & \emph{La programmation du Z-80}, R. Zaks \\
Eyrolles & \emph{L'assembleur facile du Z-80} \\
P.S.I. & \emph{Programmer en assembleur}, A. Pinaud \\
Canon & Manuel de la carte moniteur XP-140F \\
Club C7 & Gazettes \emph{Le Son du Canon} ; \emph{Les mystères du X-07} \\
Tests & \emph{L'Ordinateur Individuel} \\
S.P.E. & \emph{Micro Systèmes} \\
\bottomrule
\end{tabular}
\end{center}

\end{document}
